\graphicspath{{notes/2inner/asy/}}
\def\at#1#2{#1_{#2}}

\subsection{Orthogonal Projections}

Recall the discussion of the Gram-Schmidt process, where we saw that any finite-dimensional subspace $W$ of an inner product space $V$ has an orthonormal basis $\beta_W=\{\vect{\vw}{n}\}$. In such a situation, we can define the \emph{orthogonal projections} onto $W$ and $W^\perp$ via
\[\pi_W:V\to W:\vx\mapsto \sum_{j=1}^n\ip{\vx,\vw_j}\vw_j,\qquad \pi_W^\perp:V\to W^\perp:\vx\mapsto \vx-\pi_W(\vx)\]
Our previous goal was to use orthonormal bases to ease \emph{computation.} In this section we develop projections more generally. First recall the notion of a direct sum within a vector space $V$:
\[V=X\oplus Y \iff \forall \vv\in V,\ \exists \text{ unique }\vx\in X,\ \vy\in Y\text{ such that }\vv=\vx+\vy\]

\begin{defn}[lower separated=false, sidebyside, sidebyside align=top seam, sidebyside gap=0pt, righthand width=0.25\linewidth]{}{}
A linear map $\rT\in\cL(V)$ is a \emph{projection} if:
\[V=\cR(\rT)\oplus\cN(\rT)\quad\text{and}\quad \at\rT{\cR(\rT)}=\at\I{\cR(\rT)}\]
Otherwise said, $\rT(\vr+\vn)=\vr$ whenever $\vr\in\cR(\rT)$ and $\vn\in\cN(\rT)$.\medbreak
Alternatively, given $V=X\oplus Y$, the \emph{projection along $Y$ onto $X$} is the map $\vv=\vx+\vy\mapsto\vx$.\medbreak
We call a $A\in M_n(\F)$ a \emph{projection matrix} if $\rL_A\in\cL(\F^n)$ is a projection.
\tcblower
\flushright\includegraphics{orth-nonproj}
\end{defn}

\begin{example}{}{proj}
$A=\frac 15\begin{smatrix}
6&-2\\3&-1
\end{smatrix}$ is a projection matrix with $\cR(A)=\Span\stwovec 21$ and $\cN(A)=\Span\stwovec 13$.\smallbreak
Indeed, it is straightforward to describe all projection matrices in $M_2(\R)$. There are three cases:
\begin{enumerate}
  \item $A=I$ is the identity matrix: $\cR(A)=\R^2$ and $\cN(A)=\{\V0\}$;
  \item $A=0$ is the zero matrix: $\cR(A)=\{\V0\}$ and $\cN(A)=\R^2$;
  \item Choose \emph{distinct} subspaces $\cR(A)=\Span\stwovec ab$ and $\cN(A)=\Span\stwovec cd$, then
  \[A=\frac 1{ad-bc}\twovec ab(d\ -c)=\frac 1{ad-bc}\begin{pmatrix}
  ad&-ac\\bd&-bc
  \end{pmatrix}\]
  Think about why this last does what we claim.
\end{enumerate}
\end{example}

It should be clear that every projection $\rT$ has (at most) two eigenspaces:
\begin{itemize}
  \item $\cR(\rT)$ is an eigenspace with eigenvalue 1
  \item $\cN(\rT)$ is an eigenspace with eigenvalue 0
\end{itemize}
If $V$ is finite-dimensional and $\rho,\eta$ are bases of $\cR(T)$, $\cN(T)$ respectively, then the matrix of $\rT$ with respect to $\rho\cup\eta$ has block form
\[[\rT]_{\rho\cup\eta} =\left(\begin{array}{c|c}
I&0\\\hline
0&0
\end{array}\right)\]
where $\Rank I=\Rank\rT$. In particular, every finite-dimensional projection is diagonalizable.
\goodbreak

\begin{lemm}{}{projnotorth}
$\rT\in\cL(V)$ is a projection if and only if $\rT^2=\rT$.
\end{lemm}

\begin{proof}
Throughout, assume $\vr\in\cR(\rT)$ and $\vn\in\cN(\rT)$.\smallbreak
$(\Rightarrow)$\quad Since every vector in $V$ has a unique representation $\vv=\vr+\vn$, simply compute
\[\rT^2(\vv)=\rT\bigl(\rT(\vr+\vn)\bigr)=\rT(\vr)=\vr=\rT(\vv)\]
$(\Leftarrow)$\quad Suppose $\rT^2=\rT$. Note first that if $\vr\in\cR(\rT)$, then $\vr=\rT(\vv)$ for some $\vv\in V$, whence
\[\rT(\vr)=\rT^2(\vv)=\rT(\vv)=\vr\tag{$\dag$}\]
Thus $\rT$ is the identity on $\cR(\rT)$. Moreover, if $\vx\in\cR(\rT)\cap\cN(\rT)$, ($\dag$) says that $\vx=\rT(\vx)=\V0$, whence
\[\cR(\rT)\cap\cN(\rT)=\{\V0\}\]
and so $\cR(\rT)\oplus\cN(\rT)$ is a well-defined \emph{subspace} of $V$.\footnotemark To finish things off, let $\vv\in V$ and observe that
\[\rT\bigl(\vv-\rT(\vv)\bigr)=\rT(\vv)-\rT^2(\vv)=\V0\implies \vv-\rT(\vv)\in\cN(\rT)\]
so that $\vv=\rT(\vv)+\bigl(\vv-\rT(\vv)\bigr)$ is a decomposition into $\cR(\rT)$- and $\cN(\rT)$-parts. We conclude that $V=\cR(\rT)\oplus\cN(\rT)$ and that $\rT$ is a projection.
\end{proof}

\footnotetext{In finite dimensions, the rank--nullity theorem and dimension counting finishes the proof here without having to proceed further:
\[\dim(\cR(\rT)\oplus \cN(\rT))=\Rank\rT+\Null\rT=\dim V\implies \cR(\rT)\oplus \cN(\rT)=V\]}

\goodbreak




Thus far the discussion hasn't had anything to do with inner products\ldots

\begin{defn}{}{}
An \emph{orthogonal projection} is a projection $\rT\in\cL(V)$ on an inner product space for which we additionally have
\[\cN(\rT)=\cR(\rT)^\perp\quad\text{and}\quad \cR(\rT)=\cN(\rT)^\perp\]
Alternatively, given $V=W\oplus W^\perp$, the \emph{orthogonal projection $\pi_W$} is the projection along $W^\perp$ onto $W$: that is
\[\cR(\pi_W)=W\quad\text{and}\quad \cN(\pi_W)=W^\perp\]
The complementary orthogonal projection $\pi_W^\perp=\I-\pi_W$ has $\cR(\pi_W^\perp)=W^\perp$ and $\cN(\pi_W^\perp)=W$.
\end{defn}

\begin{example*}{\ref{ex:proj} continued}{}
The identity and zero matrices are both $2\times 2$ orthogonal projection matrices, while those of type 3 are orthogonal if $\stwovec ab\cdot\stwovec cd=0$: we obtain
\[A=\frac 1{a^2+b^2}\twovec ab(a\ b)=\frac 1{a^2+b^2}\begin{pmatrix}
a^2&ab\\ab&b^2
\end{pmatrix}\]
More generally, if $W\le\F^n$ has orthonormal basis $\{\vect{\vw}{k}\}$, then the matrix of $\pi_W$ is $\sum_{j=1}^k\vw_j\vw_j^*$.
\end{example*}

\goodbreak

\begin{thm}{}{}
A projection $\rT\in\cL(V)$ is orthogonal if and only if it is self-adjoint $\rT=\rT^*$. 
\end{thm}
  
\begin{proof}
($\Rightarrow$)\quad By assumption, $\cR(\rT)$ and $\cN(\rT)$ are orthogonal subspaces. Letting $\vx,\vy\in V$ and using subscripts to denote $\cR(\rT)$- and $\cN(\rT)$-parts, we see that
\[\ip{\vx,\rT(\vy)}=\ip{\vx_r+\vx_n,\vy_r}=\ip{\vx_r,\vy_r+\vy_n}=\ip{\rT(\vx),\vy} \implies \rT^*=\rT\]
($\Leftarrow$)\quad Suppose $\rT$ is a self-adjoint projection. By the fundamental subspaces theorem,
\[\cN(\rT)=\cN(\rT^*)=\cR(\rT)^\perp\]
Since $\rT$ is a projection already, we have $V=\cR(\rT)\oplus\cN(\rT)= \cR(\rT)\oplus \cR(\rT)^\perp$, from which 
\[\cR(\rT)=\bigl(\cR(\rT)^\perp\bigr)^\perp=\cN(\rT)^\perp\tag*{\qedhere}\]
\end{proof}
 

The language of projections allows us to rephrase the Spectral Theorem.
  
\begin{thm}{Spectral Theorem, mk.\,II}{}
Let $V$ be a finite-dimensional complex/real inner product space and $\rT\in\cL(V)$ be normal/self-adjoint with spectrum $\{\lambda_1,\ldots,\lambda_k\}$ and corresponding eigenspaces $E_1,\ldots,E_k$. Let $\pi_j\in\cL(V)$ be the orthogonal projection onto $E_j$. Then:
\begin{enumerate}
  \item $V=E_1\oplus \cdots\oplus E_k$ is a direct sum of orthogonal subspaces, in particular, $E_j^\perp$ is the direct sum of the remaining eigenspaces.
  \item $\pi_i\pi_j=0$ if $i\neq j$.
  \item (Resolution of the identity)\quad $\I_V=\pi_1+\cdots+\pi_k$
  \item (Spectral decomposition)\quad $\rT=\lambda_1\pi_1+\cdots+\lambda_k\pi_k$
\end{enumerate}
\end{thm}

\begin{proof}
\begin{enumerate}
  \item $\rT$ is diagonalizable and so $V$ is the direct sum of the eigenspaces of $\rT$. Since $\rT$ is normal, the eigenvectors corresponding to distinct eigenvalues are orthogonal, whence the eigenspaces are mutually orthogonal. In particular, this says that
  \[\hat E_j:=\bigoplus_{i\neq j}E_i\le E_j^\perp\]
  Since $V$ is finite-dimensional, we have $V=E_j\oplus E_j^\perp$, whence
  \[\dim \hat{E_j}=\sum\limits_{i\neq j}\dim E_i=\dim V-\dim E_j=\dim E_j^\perp \implies \hat E_j=E_j^\perp\]
  \item This is clear by part 1, since $\cN(\pi_j)=E_j^\perp=\hat E_j$.
  \item Write $\vx=\sum_{j=1}^k\vx_j$ where each $\vx_j\in E_j$. Then $\pi_j(\vx)=\vx_j$: now add\ldots
  \item $\rT(\vx)=\sum_{j=1}^k\rT(\vx_j)=\sum_{j=1}^k\lambda_j\vx_j=\sum_{j=1}^k\lambda_j\pi_j(\vx)$.\qedhere
\end{enumerate}
\end{proof}

\goodbreak


\begin{examples}{}{}
We verify the resolution of the identity and the spectral decomposition; for clarity, we index projections and eigenspaces by eigenvalue rather than the natural numbers.
\begin{enumerate}
  \item The symmetric matrix $A=\begin{smatrix}
  10&2\\2&7
  \end{smatrix}$ has spectrum $\{6,11\}$ and orthonormal eigenvectors
  \[\vw_6=\frac 1{\sqrt 5}\twovec 1{-2},\qquad \vw_{11}=\frac 1{\sqrt 5}\twovec 21\]
  The corresponding projections therefore have matrices
  \[\pi_6=\vw_6\vw_6^T=\frac 15\twovec 1{-2}(1\ \ -2)
  =\frac 15\begin{pmatrix}
  1&-2\\-2&4
  \end{pmatrix},\qquad
  \pi_{11}=\vw_{11}\vw_{11}^T
  =\frac 15\begin{pmatrix}
  4&2\\2&1
  \end{pmatrix}\]
  from which the resolution of the identity and the spectral decomposition are readily verified:
  \[\pi_6+\pi_{11}=\begin{pmatrix}
1&0\\0&1
  \end{pmatrix}
  \quad\text{and}\quad
  6\pi_6+11\pi_{11}= \frac 15\begin{pmatrix}
  6+44&-12+22\\
  -12+22&24+11
  \end{pmatrix}=A\]
  
  \item The normal matrix $B=\begin{smatrix}
		1+i&1+i\\
		-1-i&1+i
  \end{smatrix}\in M_2(\C)$ has spectrum $\{2,2i\}$ and corresponding
orthonormal eigenvectors
	\[\vw_2=\frac 1{\sqrt 2}\twovec i1,\qquad \vw_{2i}=\frac 1{\sqrt 2}\twovec 1i\]
	The orthogonal projection matrices are therefore
	\[\pi_2=\vw_2\vw_2^*=\frac 12\twovec i1(-i\ \ 1)=\frac 12\begin{pmatrix}
1&i\\-i&1
	\end{pmatrix}
	\qquad \pi_{2i}=\vw_{2i}\vw_{2i}^*=\frac 12\begin{pmatrix}
	1&-i\\i&1
	\end{pmatrix}\]
	from which
	\begin{gather*}
		\pi_2+\pi_{2i}=\begin{pmatrix}
		1&0\\0&1
		\end{pmatrix}\quad\text{and}\quad
		2\pi_2+2i\pi_{2i}=\begin{pmatrix}
		1&i\\-i&1
		\end{pmatrix}+\begin{pmatrix}
		i&1\\-1&i
		\end{pmatrix}=B
	\end{gather*}

	\item The matrix $C=\begin{smatrix}
		0&1&1\\
		1&0&1\\
		1&1&0
	\end{smatrix}$ has spectrum $\{-1,2\}$, an orthonormal eigenbasis
	\[\{\vu,\vv,\vw\} =\left\{\frac 1{\sqrt 3}\threevec 111, \frac 1{\sqrt 2}\threevec 1{-1}0,\frac 1{\sqrt 6}\threevec 11{-2} \right\}\]
	and eigenspaces $E_2=\Span\{\vu\}$ and $E_{-1}=\Span\{\vv,\vw\}$. The orthogonal projections have matrices
	\begin{gather*}
		\pi_2=\vu\vu^T=\frac 13\threevec 111(1\ \ 1\ \ 1)
		=\frac 13\begin{pmatrix}
			1&1&1\\
			1&1&1\\
			1&1&1
		\end{pmatrix}\\
		\begin{aligned}
			\pi_{-1}&=\vv\vv^T+\vw\vw^T
			=\frac 12\threevec 1{-1}0(1\ \ -1\ \ 0)+\frac 16\threevec 11{-2}(1\ \ 1\ \ -2)\\
			&=\frac 12\begin{pmatrix}
				1&-1&0\\
				-1&1&0\\
				0&0&0
			\end{pmatrix}+\frac 16\begin{pmatrix}
				1&1&-2\\
				1&1&-2\\
				-2&-2&4
			\end{pmatrix}
			=\frac 13\begin{pmatrix}
				2&-1&-1\\
				-1&2&-1\\
				-1&-1&2
			\end{pmatrix}
		\end{aligned}
	\end{gather*}
	It is now easy to check the resolution of the identity and the spectral decomposition:
	\[\pi_2+\pi_{-1}=I\quad\text{and}\quad 2\pi_2-\pi_{-1}=C\]
\end{enumerate}
\end{examples}

\vfil\goodbreak

\boldsubsubsection{Orthogonal Projections and Minimization Problems}\phantomsection\label{sec:leastsq}

We finish this section with an important observation that drives much of the application of inner product spaces to other parts of mathematics and beyond. Throughout this discussion, $X$ and $Y$ denote inner product spaces.


\begin{thm}{}{minthm}
Suppose $Y=W\oplus W^\perp$. For any $\vy\in Y$, the orthogonal projection $\pi_W(\vy)$ is the unique element of $W$ which minimizes the distance to $\vy$:
\[\forall\vw\in W,\ \Nm{\vy-\pi_W(\vy)}\le\Nm{\vy-\vw}\]
\end{thm}

\begin{proof}
Apply Pythagoras' Theorem: since $\pi^\perp_W(\vy)=\vy-\pi_W(\vy)\in W^\perp$,\smallbreak

\begin{minipage}[t]{0.7\linewidth}\vspace{-15pt}
	\begin{align*}
	\Nm{\vy-\vw}^2&=\Nm{\vy-\pi_W(\vy)+\pi_W(\vy)-\vw}^2\\
	&=\Nm{\vy-\pi_W(\vy)}^2+\Nm{\pi_W(\vy)-\vw}^2\\
	&\ge \Nm{\vy-\pi_W(\vy)}^2
	\end{align*}
with equality if and only if $\vw =\pi_W(\vy)$.
	\end{minipage}
	\hfill\begin{minipage}[t]{0.29\linewidth}\vspace{-30pt}
	\flushright\href{http://www.math.uci.edu/~ndonalds/math121b/orth-orthproj4.html}{\includegraphics{orth-orthproj4}}
	\end{minipage}\smallbreak
\end{proof}

This set up can be used to compute accurate approximations in many contexts.

\begin{examples}{}{minproj}
\exstart To obtain a quadratic polynomial approximation $p(x)=a+bx+cx^2$ to $e^x$ on the interval $[-1,1]$ we choose to minimize the integral $\int_{-1}^1\nm{e^x-p(x)}^2\,\dx$, namely the squared $L^2$-norm $\Nm{e^x-p(x)}^2$ on $C[-1,1]$. If we let $W=\Span\{1,x,x^2\}$, then the finite-dimensionality of $W$ means that $C[-1,1]=W\oplus W^\perp$. By the Theorem, the solution is $p(x)=\pi_W(e^x)$.
\begin{enumerate}\setcounter{enumi}{1}
    \item[]	To compute this, recall that we have an orthonormal basis for $W$, namely\smallbreak
	
	\begin{minipage}[t]{0.65\linewidth}\vspace{-15pt}
	\[\left\{\tfrac 1{\sqrt 2},\sqrt{\tfrac 32}x,\sqrt{\tfrac 58}(3x^2-1)\right\}\]
		from which
	\begin{align*}
	\textcolor{Green}{p(x)}&=\frac 12\ip{1,e^x}+\frac 32\ip{x,e^x}x+\frac 58(3x^2-1)\ip{3x^2-1,e^x}\\
	&=\frac 12\int_{-1}^1e^x\,\dx +\frac 32x\int_{-1}^1xe^x\,\dx\\
	&\qquad\qquad\qquad+\frac 58(3x^2-1)\int_{-1}^1(3x^2-1)e^x\,\dx\\
	&=\frac 12(e-e^{-1}) +3e^{-1}x +\frac 54(e-7e^{-1})(3x^2-1)\\
	&\approx \textcolor{purple}{1.18+1.10x}+0.179(3x^2-1)\\
	&\approx \textcolor{Green}{1+1.1x+0.537x^2}
	\end{align*}
	\end{minipage}\begin{minipage}[t]{0.35\linewidth}\vspace{0pt}
	\flushright\includegraphics[scale=0.9]{orth-expapprox}
	\end{minipage}\smallbreak	
	The \textcolor{purple}{linear} and \textcolor{Green}{quadratic} approximations to $\textcolor{blue}{y=e^x}$ are drawn. Compare this with the Maclaurin polynomial $e^x\approx 1+x+\frac 12x^2$ from calculus.

	
	\goodbreak
	\item\label{ex:minproj2} The $n\th$ Fourier approximation of a function $f(x)$ is its orthogonal projection onto the finite-dimensional subspace
	\begin{align*}
	W_n&=\Span\{1,e^{ix},e^{-ix},\ldots,e^{inx},e^{-inx}\}\\
	&=\Span\{1,\cos x,\sin x,\cos 2x,\sin 2x,\ldots,\cos nx,\sin nx\}
	\end{align*}
	within $L^2[-\pi,\pi]$, namely
	\begin{align*}
	\mathcal F_n(x)&=\frac 1{2\pi}\sum_{k=-n}^n\ip{f(x),e^{ikx}}e^{ikx}\\
	 &=\frac 1{2\pi}\ip{f(x),1}+\frac 1{\pi}\sum_{k=1}^n\ip{f(x),\cos kx}\cos kx+\ip{f(x),\sin kx}\sin kx
	\end{align*}
	According to the Theorem, this is the unique function in $\mathcal F_n(x)\in W_n$ minimizing the integral
	\[\Nm{f(x)-\mathcal F_n(x)}^2=\int_{-\pi}^\pi\nm{f(x)-\mathcal F_n(x)}^2\,\dx\]
	For example, if $f(x)=\begin{cases}
	1&\text{if }0<x\le\pi\\
	-1&\text{if }-\pi<x\le 0
	\end{cases}$ \space is extended periodically, then
	\[\mathcal F_{2n-1}(x)=\frac 4\pi\sum_{j=1}^n\frac{\sin(2j-1)x}{2j-1}=\frac 4\pi\left(\sin x+\frac{\sin 3x}3+\frac{\sin 5x}5+\cdots+\frac{\sin(2n-1)x}{2n-1} \right)\]
	\begin{center}
	\href{http://www.math.uci.edu/~ndonalds/math121b/orth-fourier.html}{\includegraphics{orth-fourier2}}\\
	\textcolor{blue}{$y=f(x)$} and its eleventh Fourier approximation \textcolor{Green}{$y=\mathcal F_{11}(x)$}
	\end{center}
  
%   \item To obtain a quadratic polynomial approximation $p(x)$ to $e^x$ on the interval $[0,1]$ we might wish to minimize the integral $\int_0^1\nm{e^x-p(x)}^2\,\dx$.\smallbreak
%   
%   The integral is the norm-squared $\Nm{e^x-p(x)}^2$ in the space $C[0,1]$ equipped with the $L^2$ inner product! If we let $W=\Span\{1,x,x^2\}$, then the finite-dimensionality of $W$ means that $C[0,1]=W\oplus W^\perp$. By the theorem, the solution is that $p(x)=\pi_W(e^x)$.\smallbreak
% 	To compute this, recall that we have an orthonormal basis for $W$
% 	\[\beta=\left\{1,\sqrt 3(2x-1),\sqrt{5}(6x^2-6x+1)\right\}\]
% 	from which
% 	\begin{align*}
% 	p(x)&=\ip{1,e^x}+3\ip{2x-1,e^x}(2x-1)+5\ip{6x^2-6x+1,e^x}(6x^2-6x+1)\\
% 	&=\int_0^1e^x\,\dx +3(2x-1)\int_0^1(2x-1)e^x\,\dx+5(6x^2-6x+1)\int_0^1(6x^2-6x+1)e^x\,\dx\\
% 	&=e-1 +3(3-e)(2x-1) +5(7e-19)(6x^2-6x+1)\\
% 	&\approx 1.72+0.84(2x-1)+0.2(6x^2-6x+1)=1.2x^2+0.48x+1.08
% 	\end{align*}
\end{enumerate}
\end{examples}

We'll return to related minimization problems in the next section.

\vfil

\begin{exercises}
\exstart Compute the matrices of the orthogonal projections onto $W$ viewed as subspaces of the standard inner product spaces $\R^n$ or $\C^n$.

\begin{enumerate}\setcounter{enumi}{1}
  \item[]\begin{enumerate}
    \item $W=\Span\twovec 4{-1}$\hfill \text{(b) }\space $W=\Span\left\{\threevec 121,\threevec 10{-1}\right\}$ \hfill \text{(c) }\space
    $W=\Span\left\{\threevec i10,\threevec 1i1\right\}$
    \item[(d)] $W=\Span\left\{\threevec 110,\threevec 121\right\}$\quad (\emph{watch out, these vectors aren't orthogonal!})
	\end{enumerate}
	\goodbreak

	\item For each of the following matrices, compute the projections onto each eigenspace, verify the resolution of the identity and the spectral decomposition.\vspace{-5pt}
	\[\text{(a)}\ \ \begin{pmatrix}
    1&2\\2&1
    \end{pmatrix}
    \qquad
    \text{(b)}\ \ \begin{pmatrix}
    0&-1\\1&0
    \end{pmatrix}
    \qquad
    \text{(c)}\ \ \begin{pmatrix}
    2&3-3i\\3+3i&5
    \end{pmatrix}
    \qquad
    \text{(d)}\ \ \begin{pmatrix}
    2&1&1\\1&2&1\\1&1&2
    \end{pmatrix}
    \]
    \vspace{-5pt}
    (\emph{You should have orthonormal eigenbases from the previous section})
	
	\item If $W$ be a finite-dimensional subspace of an inner product space $V$. If $\rT=\pi_W$ is the orthogonal projection onto $W$, prove that $\I-\rT$ is the orthogonal projection onto $W^\perp$.
	
	\item Let $\rT\in\cL(V)$ where $V$ is finite-dimensional.
	\begin{enumerate}
	  \item If $\rT$ is an orthogonal projection, prove that $\Nm{\rT(\vx)}\le\Nm{\vx}$ for all $\vx\in V$.
	 	\item Give an example of a projection for which the inequality in (a) is \emph{false.}
	 	\item If $\rT$ is a projection for which $\Nm{\rT(\vx)}=\Nm{\vx}$ for all $\vx\in V$, what is $\rT$?
	 	\item If $\rT$ is a projection for which $\Nm{\rT(\vx)}\le\Nm{\vx}$ for all $\vx\in V$, prove that $\rT$ is an \emph{orthogonal} projection.
	\end{enumerate}
	
	\item Let $\rT$ be a normal operator on a finite-dimensional inner product space. If $\rT$ is a projection, prove that it must be an orthogonal projection.
	
	\item\label{exs:lastqn} Let $\rT$ be a normal operator on a finite-dimensional complex inner product space $V$. Use the spectral decomposition $\rT=\lambda_1\pi_1+\cdots+\lambda_k\pi_k$ to prove:
	\begin{enumerate}
	  \item If $\rT^n$ is the zero map for some $n\in\N$, then $\rT$ is the zero map.
	  \item $\rU\in\cL(V)$ commutes with $\rT$ if and only if $\rU$ commutes with each $\pi_j$.
	  \item There exists a normal $\rU\in\cL(V)$ such that $\rU^2=\rT$.
	  \item $\rT$ is invertible if and only if $\lambda_j\neq 0$ for all $j$.
	  \item $\rT$ is a projection if and only if every $\lambda_j=0$ or 1.
	  \item $\rT=-\rT^*$ if and only if every $\lambda_j$ is imaginary.
	\end{enumerate}
	
	\item Find a linear approximation to $f(x)=e^x$ on $[0,1]$ using the $L^2$ inner product.
	
	\item Consider the $L^2$ inner product on $C[-\pi,\pi]$ inner product.
	\begin{enumerate}
	  \item Explain why $\ip{\sin x,x^{2n}}=0$ for all $n$.
	  \item Find linear and \emph{cubic} approximations to $f(x)=\sin x$.\smallbreak
	  (\emph{Feel free to use a computer algebra package to evaluate the integrals}!)
	\end{enumerate}
	
	
	\item Revisit Example \hyperref[ex:minproj2]{\ref*{ex:minproj}.\ref*{ex:minproj2}}
	\begin{enumerate}
	  \item Verify that the general complex ($e^{ikx}$) and real ($\cos kx,\sin kx$) expressions for the Fourier approximation are correct.\smallbreak
	  (\emph{Hint: use Euler's formula $e^{ikx}=\cos kx+i\sin kx$})
	  \item Verify the explicit expression for $\mathcal F_{2n-1}(x)$ when $f(x)$ is the given step-function. What is $\mathcal F_{2n}(x)$ in this case?
% 	  \item Find the $n\th$ Fourier approximation using the interval $[-T,T]$ and the subspace
% 	  \[W_n=\Span\left\{e^{\frac{i\pi kx}T}:-n\le k\le n\right\}=\Span\left\{\cos\frac{k\pi x}T,\sin\frac{k\pi x}T:0\le k\le n\right\}\]
	\end{enumerate}
\end{enumerate}
\end{exercises}
	
	\clearpage


\subsection{The Singular Value Decomposition and the Pseudoinverse}

Given $\rT\in\cL(V,W)$ between finite-dimensional inner product spaces, the overarching concern of this chapter is the existence and computation of bases $\beta,\gamma$ of $V,W$ with two properties:
\begin{itemize}
  \item That $\beta,\gamma$ be \emph{orthonormal,} thus facilitating easy calculation within $V,W$;
  \item That the matrix $[\rT]_\beta^\gamma$ be as simple as possible.
\end{itemize}

We have already addressed two special cases:
\begin{description}
	\item[Spectral Theorem] When $V=W$ and $\rT$ is normal/self-adjoint, $\exists\beta=\gamma$ such that $[\rT]_\beta$ is diagonal.
	\item[Schur's Lemma] When $V=W$ and $p(t)$ splits, $\exists\beta=\gamma$ such that $[\rT]_\beta$ is upper triangular.	
\end{description}

In this section we allow $V\neq W$ and $\beta\neq\gamma$, and obtain a result that applies to any linear map between finite-dimensional inner product spaces.

\begin{example}{}{svd1}
Let $A=\begin{smatrix}
3&1\\
2&-2\\
1&3
\end{smatrix}$ and consider orthonormal bases $\beta=\{\vv_1,\vv_2\}$ of $\R^2$ and $\gamma=\{\vw_1,\vw_2,\vw_3\}$ of $\R^3$ respectively:
\[\beta=\left\{\frac 1{\sqrt 2}\twovec 11,\frac 1{\sqrt 2}\twovec{-1}1\right\},\quad \smash{\gamma=\left\{\frac 1{\sqrt 2}\threevec 101,\frac 1{\sqrt 6}\threevec{-1}{-2}1, \frac 1{\sqrt 3}\threevec 1{-1}{-1}\right\}}\]
Since $A\vv_1=4\vw_1$ and $A\vv_2=2\sqrt 3\vw_2$, whence $[\rL_A]_\beta^\gamma=\begin{smatrix}
4&0\\
0&2\sqrt 3\\
0&0
\end{smatrix}$ is almost diagonal.
\end{example}

Our main result says that such bases always exist.

\begin{thm}{Singular Value Decomposition}{}
Suppose $V,W$ are finite-dimensional inner product spaces and that $\rT\in\cL(V,W)$ has rank $r$. Then:

\begin{enumerate}
  \item There exist orthonormal bases $\beta=\{\vect{\vv}{n}\}$ of $V$ and $\gamma=\{\vect{\vw}{m}\}$ of $W$, and positive scalars $\sigma_1\ge\sigma_2\ge\cdots\ge \sigma_r$ such that
\[\rT(\vv_j)=\begin{cases}
	\sigma_j\vw_j&\text{if }j\le r\\
	\V0&\text{otherwise}
	\end{cases}
	\qquad\text{equivalently}\quad [\rT]_\beta^\gamma=\begin{pmatrix}
	\diag(\sigma_1,\ldots,\sigma_r)&O\\
	O&O
	\end{pmatrix}\]
	\item Any such $\beta$ is an eigenbasis of $\rT^*\rT$, whence the scalars $\sigma_j$ are uniquely determined by $\rT$: indeed
  \[\rT^*\rT(\vv_j)=\begin{cases}
  \sigma_j^2\vv_j&\text{if }j\le r\\
  \V0&\text{otherwise}
  \end{cases}
  \qquad\text{and}\qquad
  \rT^*(\vw_j)=\begin{cases}
  \sigma_j\vv_j&\text{if }j\le r\\
  \V0&\text{otherwise}
  \end{cases}\]
  
	\item If $A\in M_{m\times n}(\F)$ with $\Rank A=r$, then $A=P\Sigma Q^*$ where
  \[\Sigma=[\rL_A]_\beta^\gamma=\begin{pmatrix}
	\diag(\sigma_1,\ldots,\sigma_r)&O\\
	O&O
	\end{pmatrix},\quad P=(\vect{\vw}{m}),\quad Q=(\vect{\vv}{n})\]
	Since the columns of $P,Q$ are orthonormal, these matrices are \emph{unitary.}
\end{enumerate}
\end{thm}

\goodbreak

\begin{defn}{}{}
The numbers $\sigma_1,\ldots,\sigma_r$ are the \emph{singular values} of $\rT$. If $\rT$ is not maximum rank, we have additional zero singular values $\sigma_{r+1}=\cdots=\sigma_{\min(m,n)}=0$.
\end{defn}



\begin{description}
  \item[\normalfont\emph{Freedom of Choice}\quad] While the singular values are uniquely determined, there is often significant freedom regarding the bases $\beta$ and $\gamma$, particularly if any eigenspace of $\rT^*\rT$ has dimension $\ge 2$.
  
  \item[\normalfont\emph{Special Case (Spectral Theorem)}\quad] If $V=W$ and $\rT$ is normal/self-adjoint, we may \emph{choose} $\beta$ to be an eigenbasis of $\rT$, then $\sigma_j$ is the \emph{modulus} of the corresponding eigenvalue (see Exercise \ref{exs:specsingular}).
  
  \item[\normalfont\emph{Rank-one decomposition}\quad] If we write $g_j:V\to\F$ for the linear map $g_j:\vv\mapsto \ip{\vv,\vv_j}$ (recall Riesz's Theorem), then the singular value decomposition says
  \[\rT=\sum_{j=1}^r \sigma_j\vw_jg_j\quad \text{that is}\quad \rT(\vv)=\sum_{j=1}^r \sigma_j\ip{\vv,\vv_j}\vw_j\]
  thus rewriting $\rT$ as a linear combination of rank-one maps ($\vw_jg_j:V\to W$).\\
  For matrices, $g_j$ is simply multiplication by the row vector $\vv_j^*$ and we may write\footnote{Since $\beta$ is orthonormal, it is common to write $\vv_j^*$ for the map $g_j=\langle\quad,\vv_j\rangle$ in general contexts. To those familiar with the \emph{dual space} $V^*=\cL(V,\F)$, the set $\{g_1,\ldots,g_n\}=\{\vv_1^*,\ldots,\vv_n^*\}$ is the \emph{dual basis} to $\beta$. In this course $\vv_j^*$ will only ever mean the \emph{conjugate-transpose} of a \emph{column vector} in $\F^n$. This discussion is part of why physicists write inner products differently!}
	\[A=\sum_{j=1}^r\sigma_j\vw_j\vv_j^*\]
\end{description}

\begin{example*}{\ref{ex:svd1} cont}
We apply the method in the Theorem to $A=\begin{smatrix}
3&1\\
2&-2\\
1&3
\end{smatrix}$.\smallbreak
The symmetric(!) matrix $A^TA=\begin{smatrix}
	14&2\\2&14
  \end{smatrix}$ has eigenvalues $\sigma_1^2=16$, $\sigma_2^2=12$ and orthonormal eigenvectors $\vv_1=\frac 1{\sqrt 2}\stwovec 11$, $\vv_2=\frac 1{\sqrt 2}\stwovec{-1}1$. The singular values are therefore $\sigma_1=4,\sigma_2=2\sqrt 3$. Now compute
  \[\vw_1=\frac 1{\sigma_1}A\vv_1=\frac 1{4\sqrt 2}A\twovec 11=\frac 1{\sqrt 2}\threevec 101,\quad \vw_2=\frac 1{\sigma_2}A\vv_2 =\frac 1{2\sqrt 6}A\twovec{-1}1=\frac 1{\sqrt 6}\threevec{-1}{-2}1\]
  and observe that these are orthonormal. Finally choose $\vw_3=\frac 1{\sqrt 3}\sthreevec 1{-1}{-1}$ to complete the orthonormal basis $\gamma$ of $\R^3$. A singular value decomposition is therefore
  \[
  A=\begin{pmatrix}
  3&1\\2&-2\\1&3
  \end{pmatrix}=P\Sigma Q^*
  =\begin{pmatrix}
  \frac 1{\sqrt 2}&\frac{-1}{\sqrt 6}&\frac 1{\sqrt 3}\\
  0&\frac{-2}{\sqrt 6}&\frac{-1}{\sqrt 3}\\
  \frac 1{\sqrt 2}&\frac{1}{\sqrt 6}&\frac{-1}{\sqrt 3}  
  \end{pmatrix}
  \begin{pmatrix}
  4&0\\0&2\sqrt 3\\0&0
  \end{pmatrix}
  \begin{pmatrix}
  \frac 1{\sqrt 2}&\frac 1{\sqrt 2}\\\frac{-1}{\sqrt 2}&\frac 1{\sqrt 2}
  \end{pmatrix}
	\]
	By expanding the decomposition, $A$ is expressed as the sum of rank-one matrices:
	\[\sigma_1\vw_1\vv_1^T+\sigma_2\vw_2\vv_2^T =4\threevec{\frac 1{\sqrt 2}}0{\frac 1{\sqrt 2}}
	\begin{pmatrix}
	\frac 1{\sqrt 2}&\frac 1{\sqrt 2}
	\end{pmatrix}
	+
	2\sqrt 3\threevec{\frac{-1}{\sqrt 6}}{\frac{-2}{\sqrt 6}}{\frac{1}{\sqrt 6}}
	\begin{pmatrix}
	\frac{-1}{\sqrt 2}&\frac 1{\sqrt 2}
	\end{pmatrix}
  =\begin{pmatrix}
  2&2\\0&0\\2&2
  \end{pmatrix}
  +
  \begin{pmatrix}
  1&-1\\2&-2\\-1&1
  \end{pmatrix}\]
\end{example*}

\goodbreak

\begin{proof}
\begin{enumerate}
  \item Since $\rT^*\rT$ is self-adjoint, the spectral theorem says it has an orthonormal basis of eigenvectors $\beta=\{\vect{\vv}{n}\}$. If $\rT^*\rT(\vv_j)=\lambda_j\vv_j$, then
	\[\ip{\rT(\vv_j),\rT(\vv_k)}=\ip{\rT^*\rT(\vv_j),\vv_k} =\lambda_j\ip{\vv_j,\vv_k} =\begin{cases}
	\lambda_j&\text{if }j=k\\
	0&\text{if }j\neq k
	\end{cases} \tag{$\ast$}\]
	whence every eigenvalue is a non-negative real number: $\lambda_j=\Nm{\rT(\vv_j)}^2\ge 0$.\\
	Since $\Rank\rT^*\rT=\Rank\rT=r$ (Exercise \ref{exs:ranktstart}), exactly $r$ eigenvalues are non-zero; by reordering basis vectors if necessary, we may assume
	\[\lambda_1\ge\cdots\ge\lambda_r>0\]
	If $j\le r$, define $\sigma_j:=\sqrt{\lambda_j}>0$ and $\vw_j:=\frac 1{\sigma_j}\rT(\vv_j)$, then the set $\{\vect{\vw}{r}\}$ is orthonormal ($\ast$). If necessary, extend this to an orthonormal basis $\gamma$ of $W$.
	\item If orthonormal bases $\beta$ and $\gamma$ exist such that $[\rT]_\beta^\gamma=\begin{smatrix}
	\diag(\sigma_1,\ldots,\sigma_r)&O\\
	O&O
	\end{smatrix}$, then $[\rT^*]_\gamma^\beta$ is \emph{essentially the same matrix} just that its dimensions have been reversed:
	\[[\rT^*]_\gamma^\beta=\begin{pmatrix}
	\diag(\sigma_1,\ldots,\sigma_r)&O\\
	O&O
	\end{pmatrix} \implies \rT^*\rT=\begin{pmatrix}
	\diag(\sigma_1^2,\ldots,\sigma_r^2)&O\\
	O&O
	\end{pmatrix}\]
	whence $\rT^*$ and $\rT^*\rT$ are as claimed.
	
% 	\item This is straightforward given part 1.	Observe that
% 	\[\ip{\rT^*(\vw_j),\vv_k}=\ip{\vw_j,\rT(\vv_k)}=\begin{cases}
% 	\sigma_j&\text{if $j=k$ and $k\le r$}\\
% 	0&\text{otherwise}
% 	\end{cases}\]
% 	Since $\beta$ is a basis, it follows that
% 	\[
% 	\rT^*(\vw_j)=\begin{cases}
% 	\sigma_j\vv_j&\text{if $j\le r$}\\
% 	\V0&\text{otherwise}
% 	\end{cases}
% 	\]
% 	We conclude that $\rT^*\rT(\vv_j)=\sigma^2_j\vv_j$.
	\item This is merely part 1 in the context of $\rT=\rL_A\in\cL(\F^n,\F^m)$. The orthonormal bases $\beta,\gamma$ consist of \emph{column vectors} and so the (change of co-ordinate) matrices $P,Q$ are unitary.\qedhere
\end{enumerate}
\end{proof}




\begin{examples}{}{derivpseudo}
\exstart The matrix $A=\begin{smatrix}2&3\\0&2\end{smatrix}$ has $A^TA=\begin{smatrix}
  4&6\\6&13
  \end{smatrix}$ with eigenvalues $\sigma_1^2=16$ and $\sigma_2^2=1$ and orthonormal eigenbasis
	\[\beta=\left\{\frac 1{\sqrt 5}\twovec 1{2},\frac 1{\sqrt 5}\twovec{-2}1\right\}\]
	The singular values are therefore $\sigma_1=4$ and $\sigma_2=1$, from which we obtain
	\[\gamma=\left\{\frac 1{\sigma_1}A\vv_1,\frac 1{\sigma_2}A\vv_2\right\}=\left\{\frac 1{\sqrt 5}\twovec 2{1},\frac 1{\sqrt 5}\twovec{-1}2\right\}\]
	and the decomposition
	\[A=P\Sigma Q^*=
	\begin{pmatrix}
  \frac 2{\sqrt 5}&\frac{-1}{\sqrt 5}\\
  \frac 1{\sqrt 5}&\frac 2{\sqrt 5}  
  \end{pmatrix}
  \begin{pmatrix}
  4&0\\0&1
  \end{pmatrix}
  \begin{pmatrix}
  \frac 1{\sqrt 5}&\frac 2{\sqrt 5}\\
  \frac{-2}{\sqrt 5}&\frac 1{\sqrt 5}  
  \end{pmatrix}\]
  Multiplying out, we may write $A$ as a sum of rank-one matrices
  \[A=\frac 45\twovec 21\begin{pmatrix}
  1&2
  \end{pmatrix}
  +
  \frac 15\twovec{-1}2\begin{pmatrix}
  -2&1
  \end{pmatrix}
  =\frac 45\begin{pmatrix}
  2&4\\
  1&2
  \end{pmatrix}
  +
  \frac 15\begin{pmatrix}
  2&-1\\-4&2
  \end{pmatrix}\]
  
  \goodbreak
  
\begin{enumerate}\setcounter{enumi}{1}
  \item\label{ex:derivpseudo2} The decomposition can be very messy to find in non-matrix situations. Here is a classic example where we simply observe the structure directly.\smallbreak
  The $L^2$ inner product $\ip{f,g}=\int_0^1f(x)g(x)\,\dx$ on $P_2(\R)$ and $P_1(\R)$ admits orthonormal bases
\[\beta=\left\{\sqrt 5(6x^2-6x+1),\sqrt 3(2x-1),1\right\},\qquad \gamma=\left\{\sqrt 3(2x-1),1\right\}\]
Let $\rT=\diff x$ be the derivative operator. The matrix of $\rT$ is already in the required form!
\[[\rT]_\beta^\gamma=\begin{pmatrix}
2\sqrt{15}&0&0\\
0&2\sqrt 3&0
\end{pmatrix}\]
thus $\beta,\gamma$ are suitable bases and the singular values of $\rT$ are $\sigma_1=2\sqrt{15}$ and $\sigma_2=2\sqrt 3$.\smallbreak
Since $\beta,\gamma$ are orthonormal, we could have used the adjoint method to evaluate this directly,
\[[\rT^*\rT]_\beta=([\rT]_\beta^\gamma)^T[\rT]_\beta^\gamma=\begin{pmatrix}
60&0&0\\
0&12&0\\
0&0&0
\end{pmatrix}\implies \sigma_1^2=60,\ \sigma_2^2=12\]
Up to sign, $\{[\vv_1]_\beta,[\vv_2]_\beta,[\vv_3]_\beta\}$ is therefore forced to be the standard ordered basis of $\R^3$, confirming that $\beta$ was the correct basis of $P_2(\R)$ all along! %Of course, both approaches were contingent on having orthonormal bases in the first place\ldots
\end{enumerate}
\end{examples}







% \begin{example}{}{}
% We find the singular value decomposition of the matrix $A=\begin{smatrix}
%   -1&7\\3&4\\7&1
%   \end{smatrix}$.\\[5pt]
%   According to the singular value theorem, we need to compute the orthonormal eigenpairs for
%   \[A^*A=\begin{pmatrix}
%   -1&3&7\\7&4&1
%   \end{pmatrix}\begin{pmatrix}
%   -1&7\\3&4\\7&1
%   \end{pmatrix}
%   =\begin{pmatrix}
%   59&12\\12&66
%   \end{pmatrix}\]
%   When arranged in decreasing order, these are easily seen to be\
%   \[(\vv_1,\lambda_1)=\left(\frac 15\twovec 34,75\right)\quad\text{and}\quad(\vv_2,\lambda_2)=\left(\frac 15\twovec{-4}3,50\right)\]
%   It follows that the singular values are $\sigma_1=\sqrt{75}=5\sqrt 3$ and $\sigma_2=5\sqrt 2$. Now compute the other basis vectors:
%   \[\vu_1:=\frac 1{\sigma_1}A\vv_1=\frac 1{25\sqrt 3}\sthreevec{25}{25}{25}=\frac 1{\sqrt 3}\sthreevec 111,\quad \vu_2:=\frac 1{\sigma_2}A\vv_2=\frac 1{25\sqrt 2}\sthreevec{25}{0}{-25}=\frac 1{\sqrt 2}\sthreevec 10{-1}\]
%   and we take $\vu_3=\frac 1{\sqrt 6}\sthreevec 1{-2}1$ to complete the orthonormal basis $\gamma=\{\vu_1,\vu_2,\vu_3\}$. We conclude that
%   \[A=U\Sigma V=
%   \begin{pmatrix}
%   \frac 1{\sqrt 3}&\frac 1{\sqrt 2}&\frac 1{\sqrt 6}\\
%   \frac 1{\sqrt 3}&0&\frac{-2}{\sqrt 6}\\
%   \frac 1{\sqrt 3}&\frac{-1}{\sqrt 2}&\frac 1{\sqrt 6} 
%   \end{pmatrix}
%   \begin{pmatrix}
%   5\sqrt 3&0\\0&5\sqrt 2\\0&0
%   \end{pmatrix}
%   \begin{pmatrix}
%   \frac 35&\frac{-4}5\\\frac 45&\frac 35 
%   \end{pmatrix}^T\]
% \end{example}



\boldsubsubsection{The Pseudoinverse}

The singular value decomposition of a map $\rT\in\cL(V,W)$ gives rise to a natural map from $W$ back to $V$. This map behaves somewhat like an inverse even when the operator is non-invertible!

\begin{defn}{}{}
Given the singular value decomposition of a rank $r$ map $\rT\in\cL(V,W)$, the \emph{pseudoinverse} of $\rT$ is the linear map $\rT^\dag\in\cL(W,V)$ defined by
  \[\rT^\dag(\vw_j)=\begin{cases}
  \frac 1{\sigma_j}\vv_j&\text{if }j\le r\\
  \V0&\text{otherwise}
  \end{cases}\]
\end{defn}

Restricted to $\Span\{\vect vr\}\to\Span\{\vect wr\}$, the pseudoinverse really does invert $\rT$:

\begin{minipage}[t]{0.34\linewidth}\vspace{0pt}
\begin{gather*}
\rT^\dag\rT(\vv_j)=\begin{cases}
  \vv_j&\text{if }j\le r\\
  \V0&\text{otherwise}
  \end{cases}\\[10pt]
  \rT\rT^\dag(\vw_j)=\begin{cases}
   \vw_j&\text{if }j\le r\\
   \V0&\text{otherwise}
   \end{cases}
\end{gather*}
\end{minipage}\begin{minipage}[t]{0.65\linewidth}\vspace{-2pt}
\flushright$\xymatrix @C0pt @R23pt{
	V \ar@<-0.4ex>@[red][dd]_{\textcolor{red}{\rT}}
		& = 
		& \underbrace{\Span\{\vect vr\}}_{\cN(\rT)^\perp=\cR(\rT^\dag)} \ar@<-0.4ex>@[red]@{-}[d]
		& \oplus
		& \underbrace{\Span\{\vv_{r+1},\ldots,\vv_n\}}_{\cN(\rT)=\cR(\rT^\dag)^\perp} \ar@[red][ddrr]
		&& \{\V0_V\} 
		\\
&&\text{bijection} \ar@<-0.4ex>@[blue][u] \ar@<-0.4ex>@[red][d]&&&\\
	W \ar@<-0.4ex>@[blue][uu]_{\textcolor{blue}{\rT^\dag}}
		& =
		& \overbrace{\Span\{\vect wr\}}^{\cR(\rT)=\cN(\rT^\dag)^\perp} \ar@<-0.4ex>@[blue]@{-}[u]
		& \oplus
		& \overbrace{\Span\{\vw_{r+1},\ldots,\vw_m\}}^{\cR(\rT)^\perp=\cN(\rT^\dag)} \ar@[blue][uurr]
		&& \{\V0_W\}
	}
$
\end{minipage}\smallbreak
Otherwise said, the combinations are \emph{orthogonal projections}: $\rT^\dag\rT=\pi^\perp_{\cN(\rT)}$ and $\rT\rT^\dag=\pi_{\cR(\rT)}$\smallbreak

\pagebreak
Given the singular value decomposition of a matrix $A=P\Sigma Q^*$, its pseudoinverse is that of matrix of $(\rL_A)^\dag$, namely%\footnote{This \emph{isn't} the singular value decomposition of $A^\dag$ since the values $\sigma_1^{-1}\le\sigma_2^{-1}\le\cdots\le\sigma_r^{-1}$ are in the wrong order!}
\[A^\dag=\sum_{j=1}^r\frac 1{\sigma_j}\vv_j\vw_j^*=Q\Sigma^\dag P^* \quad\text{where}\quad\Sigma^\dag=\begin{pmatrix}
	\diag(\sigma_1^{-1},\ldots,\sigma_r^{-1})&O\\
	O&O
	\end{pmatrix}\]
	
% 	Restricted to $\cN(\rT)^\perp=\Span\{\vect{\vv}{r}\}$ we obtain an isomorphism\\[-5pt]
% 	\begin{minipage}[t]{0.55\linewidth}\vspace{-8pt}
% 	\[\at{\rT}{\cN(\rT)^\perp}:\cN(\rT)^\perp\to\cR(\rT)=\Span\{\vect{\vw}{r}\} \]
% 	with inverse $\at{\rT^\dag}{\cR(\rT)}:\cR(\rT)\to\cN(\rT)^\perp$.
% 	Moreover, in terms of orthogonal projections,
% 	\[\rT\rT^\dag=\pi_{\cR(\rT)}\quad\text{and}\quad \rT^\dag \rT=\pi^\perp_{\cN(\rT)}\]
% 	\end{minipage}\begin{minipage}[t]{0.45\linewidth}\vspace{0pt}
% 	\flushright\includegraphics{pseudo-commd}
% 	\end{minipage}

\begin{examples}{}{}
\exstart Again continuing Example \ref{ex:svd1}, $A=\begin{smatrix}
  3&1\\2&-2\\1&3
\end{smatrix}$ has pseudoinverse
  \begin{align*}
  A^\dag&=\frac 1{\sigma_1}\vv_1\vw_1^*+\frac 1{\sigma_2}\vv_2\vw_2^*\\
  &=\frac 1{4\sqrt 2\sqrt 2}\twovec 11(1\ 0\ 1) +\frac 1{2\sqrt 3\sqrt 2\sqrt 6}\twovec{-1}1(-1\ -2\ 1)\\
  &=\frac 18\begin{pmatrix}
  1&0&1\\1&0&1
  \end{pmatrix}+\frac 1{12}\begin{pmatrix}
  1&2&-1\\-1&-2&1
  \end{pmatrix}
  =\frac 1{24}
  \begin{pmatrix}
  5&4&1\\1&-4&5
  \end{pmatrix}
  \end{align*}
  which is exactly what we would have found by computing $A^\dag=Q\Sigma^\dag P^*$. Observe that
  \[A^\dag A=
  \begin{pmatrix}
  1&0\\0&1
  \end{pmatrix}\quad\text{and}\quad
  AA^\dag=\frac 13
  \begin{pmatrix}
  2&1&1\\1&2&1\\1&1&2
  \end{pmatrix}\]
  are the orthogonal projection matrices onto the spaces $\cN(A)^\perp=\Span\{\vv_1,\vv_2\}=\R^2$ and $\cR(A)=\Span\{\vw_1,\vw_2\}\le\R^3$ respectively. Both spaces have dimension 2, since $\Rank A=2$.
%    It also easy to check that
%   \[A(A^TA)^{-1}A^T=AA^\dag\]
% 	in accordance with our discussion of least-squares problems.
	
\begin{enumerate}\setcounter{enumi}{1}	
	\item %Finally we check that the pseudoinverse of the derivative behaves roughly as expected.\smallbreak
The pseudoinverse of $\rT=\diff x:P_2(\R)\to P_1(\R)$, as seen in Example \hyperref[ex:derivpseudo2]{\ref*{ex:derivpseudo}.\ref*{ex:derivpseudo2}}, maps
\begin{gather*}
\rT^\dag(\sqrt 3(2x-1))=\frac 1{2\sqrt{15}}\sqrt 5(6x^2-6x+1) =\frac 1{2\sqrt 3}(6x^2-6x+1)\\
\rT^\dag(1)=\frac 1{2\sqrt 3}\sqrt 3(2x-1)=x-\frac 12\\
\begin{aligned}
\implies \rT^\dag(a+bx)&=\rT^\dag\left(a+\frac b2+\frac b{2\sqrt 3}\sqrt 3(2x-1)\right)\\
&=\left(a+\frac b2\right)\left(x-\frac 12\right)+\frac b{12}(6x^2-6x+1)\\
&=\frac b2x^2+ax-\frac a2-\frac b6
\end{aligned}
\end{gather*}
The pseudoinverse of `differentiation' therefore returns a particular choice of anti-derivative, namely the unique anti-derivative of $a+bx$ lying in $\Span\{\sqrt 5(6x^2-6x+1),\sqrt 3(2x-1)\}$.
\end{enumerate}
	
\end{examples}

% \begin{example}{}{}
% If $T=L_A$ where $A=\begin{smatrix}
%   -1&7\\3&4\\7&1
%   \end{smatrix}$, then we have pseudoinverse
%   \[T^\dag:=\frac 1{\sigma_1}\vv_1\vu_1^*+\frac 1{\sigma_2}\vv_2\vu_2^*\]
%   which has matrix
%   \[A^\dag=\frac 1{5\sqrt 3\cdot 5\cdot\sqrt 3}\twovec 34(1\ 1\ 1) +\frac 1{5\sqrt 2\cdot 5\cdot\sqrt 2}\twovec{-4}3(1\ 0\ -1) =\frac 1{150}
%   \begin{pmatrix}
%   -6&6&18\\17&8&-1
%   \end{pmatrix}\]
%   It is easy to check that
%   \[A^\dag A=\begin{pmatrix}
%   1&0\\0&1
%   \end{pmatrix},
%   \quad\text{and}\quad
%   AA^\dag=\frac 16\begin{pmatrix}
%   5&2&-1\\2&2&2\\-1&2&5
%   \end{pmatrix}\]
%   The first result is unsurprising: since $T$ has rank 2 the restriction $T:\R^2\to\cR(T)$ is a isomorphism!\\
%   Finally, observe that $TT^\dag=\vu_1\vu_1^*+\vu_2\vu_2^*=\pi_{\vu_1}+\pi_{\vu_2}$ is precisely the \emph{orthogonal projection} onto $\cR(T)=\Span\{\vu_1,\vu_2\}$: indeed, given $A$ as above, it is easy to check that
%   \[A(A^TA)^{-1}A^T=AA^\dag\]
%   is the projection matrix onto the column space of $A$, as discussed in Theorem \ref{thm:leastsq}.
%   we see that \[A(A^TA)^{-1}A^T=U\Sigma V(V^T\Sigma^TU^TU\Sigma V)^{-1}V^T\Sigma^TU^T= U\begin{smatrix}
%   1&0&0\\0&1&0\\0&0&0
%   \end{smatrix}U^T=AA^\dag\]
% \end{example}


% \begin{itemize}
%   \item : the closest vector $\vr\in\cR(T)$ to a given $\vu\in U$ is $\vr:=TT^\dag(\vu)$.
%   \item In matrix language, if $A\in M_{m\times n}(\F)$ has rank $n$, then,
%   \[A^\dag=(A^*A)^{-1}A^*\]
%   In particular, if $A\vx=\vb$ is an overdetermined system of linear equations, then $A^\dag A=I_n$ and $AA^\dag=\pi_{R(A)}$, whence $\vx=A^\dag\vb$ is the unique least-squares solution to $A\vx=\vb$: it is the unique vector minimizing the distance $\Nm{A\vx-\vb}$.
% \end{itemize}
% 
% In the above example, for instance, take $\vb=\sthreevec 100$ and consider the overdetermined system
% \[\begin{cases}
% -x+7y=1\\
% 3x+4y=0\\
% 7x+y=0
% \end{cases} \quad\leftrightsquigarrow A\vx=\vb\]
% The unique least-squares solution is
% \[\vx=A^\dag\vb=\frac 1{150}
%   \begin{pmatrix}
%   -6&6&18\\17&8&-1
%   \end{pmatrix}\threevec 100 =\frac 1{150}\twovec{-6}{17} \implies A\vx=\threevec{5/6}{2/6}{-1/6}\]

\vfil\vfil

\begin{exercises}
\exstart\label{exs:pseudo1} Find the ingredients $\beta,\gamma$ and the singular values for each of the following:
\begin{enumerate}\setcounter{enumi}{1}
  \item[]\begin{enumerate}
    \item $\rT\in\cL(\R^2,\R^3)$ where $\rT\stwovec xy=\sthreevec x{x+y}{x-y}$
    \item $\rT:P_2(\R)\to P_1(\R)$ and $\rT(f)=f''$ where $\ip{f,g}:=\int_0^1f(x)g(x)\,\dx$
    \item $V=W=\Span\{1,\sin x,\cos x\}$ and $\ip{f,g}=\int_0^{2\pi}f(x)g(x)\,\dx$, with $\rT(f)=f'+2f$
  \end{enumerate}
  
  \goodbreak
  
  \item\label{exs:pseudo2} Find a singular value decomposition of each of the matrices:
  \begin{enumerate}
    \item $\begin{pmatrix}
    1&1\\1&1\\-1&-1
    \end{pmatrix}$\qquad
    (b)\ \ $\begin{pmatrix}
    1&0&1\\1&0&-1
    \end{pmatrix}$\qquad
    (c)\ \ $\begin{pmatrix}
    1&1&1\\1&-1&0\\1&0&-1
    \end{pmatrix}$
  \end{enumerate}
  
  \item Find an explicit formula for $\rT^\dag$ for each map in Exercise \hyperref[exs:pseudo1]{1}.
  
  \item Find the pseudoinverse of each of the matrix in Exercise \ref{exs:pseudo2}.
  
  \item Suppose $A=P\Sigma Q^*$ is a singular value decomposition.
  \begin{enumerate}
    \item Describe a singular value decomposition of $A^*$.
    \item Explain why $A^\dag=Q\Sigma^\dag P^*$ \emph{isn't} a singular value decomposition of $A^\dag$; what would be the correct decomposition?   	(\emph{Hint: what is wrong with $\Sigma^\dag$?})
  \end{enumerate}
  
  \item Suppose $\rT:V\to W$ is written according to the singular value theorem. Prove that $\gamma$ is a basis of eigenvectors of $\rT\rT^*$ with the \emph{same} non-zero eigenvalues as $\rT^*\rT$, including repetitions.
  
  \item\label{exs:specsingular}\begin{enumerate}
    \item Suppose $\rT=\cL(V)$ is normal. Prove that each $\vv_j$ in the singular value theorem may be chosen to be an eigenvector of $\rT$ and that $\sigma_j$ is the \emph{modulus} of the corresponding eigenvalue.
    \item Let $A=\begin{smatrix}
    0&1\\1&0
    \end{smatrix}$. Show that \emph{any} orthonormal basis $\beta$ of $\R^2$ satisfies the singular value theorem. What is $\gamma$ here? What is it about the \emph{eigenvalues} of $A$ that make this possible?\smallbreak
    (\emph{Even when $\rT$ is self-adjoint, the vectors in $\beta$ need not also be eigenvectors of $\rT$!})
  \end{enumerate}
  
  \item\label{exs:ranktstart} In the proof of the singular value theorem we claimed that $\Rank \rT^*\rT=\Rank\rT$. Verify this by checking explicitly that $\cN(\rT^*\rT)=\cN(\rT)$.\smallbreak
  (\emph{This is circular logic if you use the decomposition, so you must do without!})
  

  
 
  \item Let $V,W$ be finite-dimensional inner product spaces and $\rT\in\cL(V,W)$. Prove:
  \begin{enumerate}
    \item $\rT^*\rT\rT^\dag=\rT^\dag\rT\rT^*=\rT^*$.\\
    (\emph{Hint: evaluate on the basis $\gamma=\{\vect wm\}$ in the singular value theorem})
    \item If $\rT$ is injective, then $\rT^*\rT$ is invertible and $\rT^\dag=(\rT^*\rT)^{-1}\rT^*$.
    \item If $\rT$ is surjective, then $\rT\rT^*$ is invertible and $\rT^\dag=\rT^*(\rT\rT^*)^{-1}$.
  \end{enumerate}
  
  \item Consider the equation $\rT(\vx)=\vb$, where $\rT$ is a linear map between finite-dimensional inner product spaces. A \emph{least-squares solution} is a vector $\vx$ which minimizes $\Nm{\rT(\vx)-\vb}$.
  \begin{enumerate}
    \item Prove that $\vx_0=\rT^\dag(\vb)$ is a least-squares solution and that any other has the form $\vx_0+\vn$ for some $\vn\in\cN(\rT)$.\\
		(\emph{Hint: Theorem \ref{thm:minthm} says that $\vx_0$ is a least-squares solution if and only if $\rT(\vx_0)=\pi_{\cR(\rT)}(\vb)$})
    \item Prove that $\vx_0=\rT^\dag(\vb)$ has smaller norm than any other least-squares solution.
		\item If $\rT$ is injective, prove that $\vx_0=\rT^\dag(\vb)$ is the unique least-squares solution.
	\end{enumerate}

	
% 	\item For the following systems, find the minimal norm solution to the first and the vector which comes closest to solving the second:
% 	\[\begin{cases}
% 		-x+7y=6\\
% 		3x+4y=7\\
% 		7x+y=8
% 		\end{cases}
% 		\qquad\qquad
% 		\begin{cases}
% 		-x+7y=1\\
% 		3x+4y=0\\
% 		7x+y=0
% 		\end{cases}\]
		
	\item Find the minimal norm solution to the first system, and the least-squares solution to the second:
	\[\begin{cases}
		3x+2y+z=9\\
		x-2y+3z=3
		\end{cases}
		\qquad\qquad
		\begin{cases}
		3x+y=1\\
		2x-2y=0\\
		x+3y=0
		\end{cases}\]
\end{enumerate}
\end{exercises}

\clearpage

% \boldsubsubsection{Linear Regression and least-squares problems (non-examinable)}
% 
% Recall Exercise 8. The vector $\vx=A^\dag\vb$ is often known as a \emph{least-squares solution}\footnote{This is something of a misnomer; least-squares solutions are typically most useful when the equation has no solutions and we want to get as close as we can to a solution.} to the linear equation $A\vx=\vb$, since to minimize the (squared-)norm $\Nm{A\vx-\vb}^2$ means minimizing a sum of squares. 
% 
% \begin{defn}{}{}
% Suppose $\rT\in\cL(X,Y)$ and $\vy\in Y$. A \emph{least-squares solution} to the linear equation $\rT(\vx)=\vy$ is any vector $\vx_0\in X$ which minimizes $\Nm{\rT(\vx)-\vy}$.
% \end{defn}
% 
% In a common special situation, there is a unique least-squares solution.
% 
% \begin{thm}{}{}
% Suppose $X,Y$ are finite-dimensional and that $\rT\in\cL(X,Y)$ is injective. Then the unique least squares solution to $\rT(\vx)=\vy$ is given by the pseudoinverse
% \[\vx_0=\rT^\dag(\vy)=(\rT^*\rT)^{-1}\rT(\vy)\]
% \end{thm}
% 
% \begin{proof}
% By Theorem \ref{thm:minthm}, $\vx_0$ is a least-squares solution if and only if
% \[\rT(\vx_0)=\pi_{\cR(\rT)}(\vy)=\rT\rT^\dag(\vy)\]
% Since $\rT$ is injective, it follows that $\vx_0=\rT^\dag(\vy)$. Exercise 7 gives the alternative expression.
% \end{proof}
% 
% Provided $X,Y$ are finite-dimensional, problem 8 shows that $\vx=\rT^\dag(\vy)$ is a least-squares solution.
% 
% If $\vy\in\cR(\rT)$, these are simply solutions: plainly $\rT(\vx_0)=\vy$ minimizes $\Nm{\vy-\rT(\vx_0)}$.\smallbreak
% Otherwise, we want $\vx_0$ to come as close as possible to solving $\rT(\vx)=\vy$. With a little simplification, this can be attacked using a general method.
% 
% \begin{itemize}
%   \item If $Y=\cR(\rT)\oplus \cR(\rT)^\perp$, the Theorem says that $\rT(\vx_0)=\pi_{\cR(\rT)}(\vy)$: equivalently
%   \[\vy-\rT(\vx_0)\in\cR(\rT)^\perp\]
%   \item Additionally, if $\rT$ has an adjoint, the Fundamental Subspaces Theorem says that $\vx_0$ is a least-squares solution if and only if
%   \[
% 	\vy-\rT(\vx_0)\in \cN(\rT^*) \iff \rT^*\rT(\vx_0)=\rT^*(\vy)
% 	\]
% \end{itemize}
% These conditions are certainly met when $X,Y$ are finite-dimensional.
% 
% \begin{thm}{}{}
% Suppose $X,Y$ are finite-dimensional and $\rT\in\cL(X,Y)$. Then:
% \begin{enumerate}
%   \item $\Rank\rT^*\rT=\Rank\rT=\Rank\rT^*$.
%   \item $\rT^*\rT$ and $\rT^*$ have the same range and so $\rT^*\rT(\vx_0)=\rT^*(\vy)$ always has at least one solution.
%   \item If $\Rank\rT=\dim X$ (equivalently $\Null\rT=0$), then $\rT(\vx)=\vy$ has a unique least-squares solution $\vx_0=(\rT^*\rT)^{-1}\rT^*(\vy)$. Moreover, the orthogonal projection onto $\cR(\rT)$ is
%   \[\pi_{\cR(\rT)}=\rT(\rT^*\rT)^{-1}\rT^*\]
% \end{enumerate}
% \end{thm}
% 
% \goodbreak
% 
% \begin{proof}
% \begin{enumerate}
%   \item $\Rank\rT^*=\Rank\rT$ holds for any maps. By the rank--nullity theorem, it therefore suffices to check that $\rT^*\rT$ and $\rT$ have the same nullspace:
% 	\begin{gather*}
% 		\rT(\vx)=\V0\implies \rT^*\rT(\vx)=\V0 \quad\text{and,}\\
% 		\rT^*\rT(\vx)=\V0\implies 0=\ip{\rT^*\rT(\vx),\vx}=\ip{\rT(\vx),\rT(\vx)} =\Nm{\rT(\vx)}^2\implies \rT(\vx)=\V0
% 	\end{gather*}
% 	\item This is immediate from part 1.
% 	\item If $\Rank \rT=\dim V$, then $\rT^*\rT$ is invertible; $\vx_0$ is a least-squares solution if and only if
% 	\[\rT^*\rT(\vx_0)=\rT^*(\vy)\iff\vx_0=(\rT^*\rT)^{-1}\rT^*(\vy)\]
% 	In such a case, $\pi_{\cR(\rT)}(\vy)=\rT(\vx_0)=\rT(\rT^*\rT)^{-1}\rT^*(\vy)$.\hfill\qedhere
% \end{enumerate}
% \end{proof}
% 
% Least-squares problems are often phrased in terms of matrices: when $\rT=\rL_A$ for some $A\in M_{m\times n}(\F)$. Equivalently, we consider a system of linear equations.

% \begin{example}{}{}
% Find the unique least-squares solution to the overdetermined system
% \[\left\{\begin{array}{r@{\hspace{4pt}}c@{\hspace{4pt}}l}
% x_1+3x_2&=&5\\
% x_2&=&1\\
% x_1+2x_2&=&1
% \end{array}\right.\leftrightsquigarrow\  A\vx=\vy\quad\text{where}\quad
% A=\begin{pmatrix}
% 1&3\\
% 0&1\\
% 1&2
% \end{pmatrix},\ \vy=\threevec 511,\ X=\R^2,\ Y=\R^3\]
% The system does not have a explicit solution. Since $\Rank A=2=\dim X$, we may apply the Theorem
% \begin{align*}
% (A^TA)^{-1}A^T&=
% \left[
% \begin{pmatrix}
% 1&0&1\\
% 3&1&2
% \end{pmatrix}\begin{pmatrix}
% 1&3\\
% 0&1\\
% 1&2
% \end{pmatrix}
% \right]^{-1}
% \begin{pmatrix}
% 1&0&1\\
% 3&1&2
% \end{pmatrix}
% =\begin{pmatrix}
% 2&5\\
% 5&14
% \end{pmatrix}^{-1}
% \begin{pmatrix}
% 1&0&1\\
% 3&1&2
% \end{pmatrix}\\
% &
% =\frac 13
% \begin{pmatrix}
% 14&-5\\
% -5&2
% \end{pmatrix}
% \begin{pmatrix}
% 1&0&1\\
% 3&1&2
% \end{pmatrix} 
% =\frac 13\begin{pmatrix}
% -1&-5&4\\
% 1&2&-1
% \end{pmatrix}\\
% \implies \vx_0&=\frac 13\begin{pmatrix}
% -1&-5&4\\
% 1&2&-1
% \end{pmatrix}\threevec 511 =\twovec{-2}{2}
% \end{align*}
% We can check this explicitly using a little multi-variable calculus: let
% \[f(x_1,x_2)=\Nm{\vy-A\vx}^2= (x_1+3x_2-5)^2+(x_2-1)^2+(x_1+2x_2-1)^2\]
% then $f$ is minimized when
% \[\partials[f]{x_1}=4x_1+10x_2-12=0=\partials[f]{x_2} =10x_1+28x_2-36 \iff \twovec{x_1}{x_2}=\twovec{-2}{2}\]
% Moreover, the orthogonal projection onto the range (column space) of $A$ is left-multiplication by the matrix
% \[\pi_{\cR(A)}=A(A^TA)^{-1}A^T =
% \frac 13\begin{pmatrix}
% 1&3\\
% 0&1\\
% 1&2
% \end{pmatrix}
% \begin{pmatrix}
% -1&-5&4\\
% 1&2&-1
% \end{pmatrix} =\frac 13\begin{pmatrix}
% 2&1&1\\
% 1&2&-1\\
% 1&-1&2
% \end{pmatrix}\]
% which may be verified directly by first finding an orthonormal basis $\left\{\frac 1{\sqrt 2}\sthreevec 101,\frac 1{\sqrt 6}\sthreevec 12{-1}\right\}$ of $\cR(A)$.
% \end{example}

\boldsubsubsection{Linear Regression (non-examinable)}

Given a data set $\{(t_j,y_j):1\le j\le m\}$, we may employ the least-squares method to find a \emph{best-fitting line} $y=c_0+c_1t$; often used to predict $y$ given a value of $t$.\smallbreak
The trick is to minimize the sum of the squares of the \textcolor{Green}{vertical deviations} of the line from the data set.
\[\sum_{j=1}^m\nm{y_j-c_0-c_1t_j}^2=\Nm{\vy-A\vx}^2\quad\text{where}\quad A=\begin{pmatrix}
t_1&1\\
\vdots&\vdots\\
t_m&1
\end{pmatrix}\quad \vx=\twovec{c_1}{c_0}\quad \vy=\threevec{y_1}{\vdots}{y_m} \]
With the indicated notation, we recognize this as a least-squares problem. Indeed if there are at least two distinct $t$-values in the data set, then $\Rank A=2$ is maximal and we have a unique best-fitting line with coefficients given by
\[\twovec{c_1}{c_0}=A^\dag\vy =(A^TA)^{-1}A^T\vy\]


\begin{example}{}{lsq}
Given the data set $\{(0,1),(1,1),(2,0),(3,2),(4,2)\}$, we compute
\begin{align*}
A=\begin{pmatrix}
0&1\\
1&1\\
2&1\\
3&1\\
4&1
\end{pmatrix},\quad \vy=\begin{pmatrix}
1\\1\\0\\2\\2
\end{pmatrix} \implies \vx_0&=\twovec{c_1}{c_0} =(A^TA)^{-1}A^T\vy=\begin{pmatrix}
30&10\\
10&5
\end{pmatrix}^{-1}
\twovec{15}{6} =\frac 3{10}\twovec{1}{2}
\end{align*}
The regression line therefore has equation $y=\frac 3{10}(t+2)$.\smallbreak
The process can be applied more generally to approximate using other functions. To find the best-fitting quadratic polynomial $y=c_0+c_1t+c_2t^2$, we'd instead work with
\begin{gather*}
A=\begin{pmatrix}
t_1^2&t_1&1\\
\vdots&\vdots&\vdots\\
t_m^2&t_m&1
\end{pmatrix}\quad \vx=\threevec{c_2}{c_1}{c_0}\quad \vy=\threevec{y_1}{\vdots}{y_m} \implies \sum_{j=1}^m\nm{y_j-c_0-c_1t-c_2t_j^2}^2 =\Nm{\vy-A\vx}^2
\end{gather*}
Provided we have at least three distinct values $t_1,t_2,t_3$, the matrix $A$ is guaranteed to have rank 3 and there will be a best-fitting least-squares quadratic: in this case
\[y=\frac 1{70}(15t^2-39t+72)\]
This curve and the best-fitting straight line are shown below.
\begin{center}
\includegraphics{least-sq-line}
\qquad
\includegraphics{least-sq-quad}
\end{center}
\end{example}

 %For instance, to find the best-fitting curve of the form $y=ae^t+b$, we would apply the method to
% \[A=\begin{pmatrix}
% e^{t_1}&1\\
% \vdots&\vdots\\
% e^{t_m}&1
% \end{pmatrix}\quad \vx=\twovec{a}{b}\quad \vy=\threevec{y_1}{\vdots}{y_m} \implies \sum_{j=1}^m\nm{y_j-ae^{t_j}-b}^2 =\Nm{\vy-A\vx}^2\]

\begin{tcolorbox}[exercisestyle,title={Optional Problems\quad }]
Use a computer to invert any $3\times 3$ matrices!
\begin{enumerate}
	\item Check the calculation for the best-fitting least-squares quadratic in Example \ref{ex:lsq}.
  
  \item Find the best-fitting least-squares linear and quadratic approximations to the data set
  \[\left\{(1,2),(3,4),(5,7),(7,9),(9,12)\right\}\]
  
  \item Suppose a data set $\{(t_j,y_j):1\le j\le m\}$ has unique regression line $y=ct+d$.
  \begin{enumerate}
    \item Show that the equations $A^TA\vx_0=A^T\vy$ can be written in matrix form
    \[\begin{pmatrix}
    \sum t^2_j&\sum t_j\\
    \sum t_j&m
    \end{pmatrix}\twovec cd =\twovec{\sum t_jy_j}{\sum y_j}\]
    \item Recover the standard expressions from statistics:
    \[c=\frac{\Cov(t,y)}{\sigma^2_t}\ \text{ and }\ d=\cl y-c\cl t\]
    where
    \begin{itemize}
      \item $\displaystyle\cl t=\frac 1m\sum_{j=1}^m t_j$ and $\displaystyle\cl y=\frac 1m\sum_{j=1}^m y_j$ are the \emph{means} (averages),
      \item $\displaystyle\sigma^2_t=\frac 1m\sum_{j=1}^m(t_j-\cl t)^2$ is the \emph{variance},
      \item $\displaystyle\Cov(t,y)=\frac 1m\sum_{j=1}^m(t_j-\cl t)(y_j-\cl y)$ is the \emph{covariance.}
    \end{itemize}
	\end{enumerate}
\end{enumerate}
\end{tcolorbox}

\clearpage

\subsection{Bilinear and Quadratic Forms}

In this section we slightly generalize the idea of an inner product. Throughout, $V$ is a vector space over a field $\F$; it need not be an inner product space and $\F$ can be any field (not just $\R$ or $\C$).

\begin{defn}{}{}
A \emph{bilinear form} $B:V\times V\to\F$ is linear in each entry: $\forall \vv,\vx,\vy\in V,\lambda\in\F$,
  \[B(\lambda\vx+\vy,\vv)=\lambda B(\vx,\vv)+B(\vy,\vv),\qquad B(\vv,\lambda\vx+\vy)=\lambda B(\vv,\vx)+B(\vv,\vy)\]
Additionally, $B$ is \emph{symmetric} if $\forall \vx,\vy\in V$,\ $B(\vx,\vy)=B(\vy,\vx)$.
\end{defn}

\begin{examples}{}{bilinear}
\exstart If $V$ is a \emph{real} inner product space, then the inner product $\ip{\ ,\ }$ is a symmetric bilinear form. Note that a \emph{complex} inner product is \emph{not bilinear!}

\begin{enumerate}\setcounter{enumi}{1}
	\item\label{ex:bilinear1} If $A\in M_n(\F)$, then $B(\vx,\vy):=\vx^TA\vy$ is a bilinear form on $\F^n$. For instance, on $\R^2$,
	\[B(\vx,\vy)=\vx^T\begin{pmatrix}
	1&2\\2&0
	\end{pmatrix}\vy =x_1y_1+2x_1y_2+2x_2y_1\]
	defines a symmetric bilinear form, though not an inner product since it isn't positive definite; for example $B(\vj,\vj)=0$.
% 	
%   \item On the vector space $V=\Z_3^2$ over the finite field $\Z_3=\{0,1,2\}$ of remainders modulo 3, define
%   \[B(\vx,\vy)=\vx^T\begin{pmatrix}
% 	2&0\\1&1
% 	\end{pmatrix}\vy=2x_1y_1+x_2y_1+x_2y_2\]
% 	For instance, $B\left(\stwovec 20,\stwovec 12\right)=1$. This isn't symmetric, since $B\left(\stwovec 12,\stwovec 20\right)=2$.
\end{enumerate}
\end{examples}

As seen above, we often make use of a matrix.

\begin{defn}{}{}
Let $B$ be a bilinear form on a finite-dimensional space with basis $\epsilon=\{\vect{\vv}{n}\}$. The \emph{matrix of $B$ with respect to $\epsilon$} is the matrix $[B]_\epsilon=A\in M_n(\F)$ with $ij^\text{th}$ entry
\[A_{ij}=B(\vv_i,\vv_j)\]
\end{defn}

Given $\vx,\vy\in V$, compute their co-ordinate vectors $[\vx]_\epsilon,[\vy]_\epsilon$ with respect to $\epsilon$, then
\[B(\vx,\vy)=[\vx]_\epsilon^TA[\vy]_\epsilon\]
The set of bilinear forms on $V$ is therefore in bijective correspondence with $M_n(\F)$. Moreover,
\[B(\vy,\vx)=[\vy]_\epsilon^TA[\vx]_\epsilon=\left([\vy]_\epsilon^TA[\vx]_\epsilon\right)^T=[\vx]_\epsilon^TA^T[\vy]_\epsilon\]
Finally, if $\beta$ is another basis of $V$, then an appeal to the change of co-ordinate matrix $Q_\epsilon^\beta$ yields
\[B(\vx,\vy)=[\vx]_\epsilon^TA[\vy]_\epsilon =(Q_\beta^\epsilon[\vx]_\beta)^TA(Q_\beta^\epsilon[\vy]_\epsilon) =[\vx]_\beta^T(Q_\beta^\epsilon)^TAQ_\beta^\epsilon[\vy]_\beta \implies [B]_\beta=(Q_\beta^\epsilon)^T[B]_\epsilon Q_\beta^\epsilon\]
To summarize:

\begin{lemm}{}{}
Let $B$ be a bilinear form on a finite-dimensional vector space.
\begin{enumerate}
  \item If $A$ is the matrix of $B$ with respect to some basis, then every other matrix of $B$ has the form $Q^TAQ$ for some invertible $Q$.
  \item $B$ is symmetric if and only if its matrix with respect to any (and all) bases is symmetric.
\end{enumerate}
\end{lemm}

Naturally, the simplest situation is when the matrix of $B$ is diagonal\ldots
\goodbreak

\begin{examples}{}{}
\exstart Example \hyperref[ex:bilinear1]{\ref*{ex:bilinear}.\ref*{ex:bilinear1}} can be written
\begin{align*}
B(\vx,\vy)&=\vx^T\begin{pmatrix}
1&2\\2&0
\end{pmatrix}\vy =x_1y_1+2x_1y_2+2x_2y_1 =(x_1+2x_2)(y_1+2y_2)-4x_2y_2\\
&=\bigl(x_1+2x_2\quad x_2\bigr)\begin{pmatrix}
1&0\\0&-4
\end{pmatrix}\twovec{y_1+2y_2}{y_2} =\vx^T\begin{pmatrix}
 	1&2\\0&1
\end{pmatrix}^T\begin{pmatrix}
1&0\\0&-4
\end{pmatrix}\begin{pmatrix}
 	1&2\\0&1
\end{pmatrix}\vy %\\
% \implies &[B]_\gamma=\begin{pmatrix}
% 1&0\\0&-4
% \end{pmatrix}\quad\text{where}\quad\gamma=\left\{\twovec 10,\twovec{-2}1\right\}
\end{align*}
If $\epsilon$ is the standard basis, then $[B]_\beta=\begin{smatrix}
1&0\\0&-4
\end{smatrix}$ where $Q_\epsilon^\beta=\begin{smatrix}
 	1&2\\0&1
\end{smatrix}$. It follows that $Q_\beta^\epsilon=\begin{smatrix}
 	1&-2\\0&1
\end{smatrix}$ from which $\beta=\left\{\stwovec 10,\stwovec{-2}1\right\}$ is a diagonalizing basis.
\begin{enumerate}\setcounter{enumi}{1}
  \item In general, we may perform a sequence of \emph{simultaneous} row and column operations to diagonalize \emph{any} symmetric $B$; we require only elementary matrices $E_{ij}^{(\lambda)}$ of type III.\footnotemark For instance:
\begin{align*}
&B(\vx,\vy)=\vx^TA\vy=\vx^T
\begin{pmatrix}
	1&-2&3\\-2&0&4\\3&4&-1
\end{pmatrix}\vy \tag{$A=[B]_\epsilon$ with respect to the standard basis}\\
&E_{21}^{(2)}AE_{12}^{(2)}=\begin{pmatrix}
	1&0&3\\0&-4&10\\3&10&-1
\end{pmatrix} \tag{add twice row 1 to row 2, columns similarly}\\
&E_{31}^{(-3)}E_{21}^{(2)}AE_{12}^{(2)}E_{13}^{(-3)}=\begin{pmatrix}
	1&0&0\\0&-4&10\\0&10&-10
\end{pmatrix} \tag{subtract 3 times row 1 from row 3, etc.}\\
&E_{23}^{(1)}E_{31}^{(-3)}E_{21}^{(2)}AE_{12}^{(2)}E_{13}^{(-3)}E_{32}^{(1)}=\begin{pmatrix}
	1&0&0\\0&6&0\\0&0&-10
\end{pmatrix} \tag{add row 3 to row 2, etc.}
\end{align*}
If $\beta$ is the diagonalizing basis, the the change of co-ordinate matrix is
\[Q_\beta^\epsilon=E_{12}^{(2)}E_{13}^{(-3)}E_{32}^{(1)}
=\begin{pmatrix}
	1&2&0\\0&1&0\\0&0&1
\end{pmatrix}
\begin{pmatrix}
	1&0&-3\\0&1&0\\0&0&1
\end{pmatrix}
\begin{pmatrix}
	1&0&0\\0&1&0\\0&1&1
\end{pmatrix}
=\begin{pmatrix}
	1&-1&-3\\0&1&0\\0&1&1
\end{pmatrix}\]
from which $\beta=\left\{\sthreevec 100,\sthreevec{-1}11,\sthreevec{-3}01\right\}$. If you're having trouble believing this, invert the change of co-ordinate matrix and check that
\[B(\vx,\vy)=(x_1-2x_2+3x_3)(y_1-2y_2+3y_3)+6x_2y_2-10(-x_2+x_3)(-y_2+y_3)\]
\end{enumerate}
\end{examples}

\footnotetext{Recall that $\smash[b]{E_{ij}^{(\lambda)}}$ is the identity matrix with an additional $\lambda$ in the $ij^\text{th}$ entry.
 \begin{itemize}\itemsep0pt\parsep0pt
   \item As a \emph{column} operation (right-multiplication), $A\mapsto A\smash{E_{ij}^{(\lambda)}}$ adds $\lambda$ times the $i^\text{th}$ column to the $j^\text{th}$.
   \item As a \emph{row} operation (left-multiplication), $A\mapsto E_{ji}^{(\lambda)}A=(E_{ij}^{(\lambda)})^TA$ adds $\lambda$ times the $i^\text{th}$ row to the $j^\text{th}$.
 \end{itemize}} 

\textcolor{red}{Warning!} If $\F=\R$ then every symmetric $B$ may be diagonalized by an orthonormal basis (for the usual dot product on $\R^n$). It is \emph{very} unlikely that our algorithm will produce such! The algorithm has two main advantages over the spectral theorem: it is typically faster and it applies to vector spaces over any field. As a disadvantage, it is highly non-unique.

\goodbreak

\begin{example}{}{diagbilinear}
We diagonalize $B(\vx,\vy)=\vx^T\begin{smatrix}1&6\\6&3\end{smatrix}\vy =x_1y_1+6x_1y_2+6x_2y_1+3x_2y_2$ in three ways.
\begin{itemize}
  \item $\begin{smatrix}1&0\\-6&1\end{smatrix}
  \begin{smatrix}1&6\\6&3\end{smatrix}
  \begin{smatrix}1&-6\\0&1\end{smatrix}
  =\begin{smatrix}1&0\\0&-33\end{smatrix}=[B]_\beta$ where $\beta=\left\{\stwovec 10,\stwovec{-6}1\right\}$. This corresponds to 
  \[B(\vx,\vy)=(x_1+6x_2)(y_1+6y_2)-33x_2y_2\]
  \item $\begin{smatrix}1&-2\\0&1\end{smatrix}
  \begin{smatrix}1&6\\6&3\end{smatrix}
  \begin{smatrix}1&0\\-2&1\end{smatrix}
  =\begin{smatrix}-11&0\\0&3\end{smatrix}=[B]_\gamma$ where $\gamma=\left\{\stwovec 1{-2},\stwovec 01\right\}$. This corresponds to 
  \[B(\vx,\vy)=-11x_1y_1+3(2x_1+x_2)(2y_1+y_2)\]
  \item If $\F=\R$, we may apply the spectral theorem to see that $[B]_\eta=\smash[t]{\begin{smatrix}2+\sqrt{37}&0\\0&2-\sqrt{37}\end{smatrix}}$ is diagonal with respect to $\eta=\left\{\stwovec 6{1+\sqrt{37}},\stwovec{-1-\sqrt{37}}6\right\}$. The expression for $B(\vx,\vy)$ in these co-ordinates is disgusting, so we omit it; that $B$ can be diagonalized orthogonally doesn't mean it should be! 
  
%   is an orthogonal eigenbasis, with respect to which $B$ is diagonal
%   \begin{gather*}
%   [B]_\eta=\frac 1{74+2\sqrt{37}}
%   \begin{smatrix}6&1+\sqrt{37}\\-1-\sqrt{37}&6\end{smatrix}
%   \begin{pmatrix}1&6\\6&3\end{pmatrix}
%   \begin{smatrix}6&-1-\sqrt{37}\\1+\sqrt{37}&6\end{smatrix}
%   =\begin{smatrix}2+\sqrt{37}&0\\0&2-\sqrt{37}\end{smatrix}%\\
%   \begin{aligned}
%   B(\vx,\vy)&=(2+\sqrt{37})\frac{(6x_1+(1+\sqrt{37})x_2)(6y_1+(1+\sqrt{37})y_2)}{74+2\sqrt{37}}\\
%   &\qquad\qquad +(2-\sqrt{37})\frac{(6x_2-(1+\sqrt{37})x_1)(6y_2-(1+\sqrt{37})y_1)}{74+2\sqrt{37}}
%   \end{aligned}
  %\end{gather*}

\end{itemize}
\end{example}

\begin{thm}{}{}
Suppose $B$ is a bilinear form of a finite-dimensional space $V$ over $\F$.
\begin{enumerate}
  \item If $B$ is diagonalizable, then it is symmetric.
  \item If $B$ is symmetric and $\F$ does not have characteristic two (see aside), then $B$ is diagonalizable. 
\end{enumerate}
\end{thm}

\begin{proof}
\begin{enumerate}
  \item If $B$ is diagonalizable, $\exists\beta$ such that $[B]_\beta$ is diagonal and thus symmetric.
  \item Suppose $B$ is non-zero (otherwise the result is trivial). We prove by induction on $n=\dim V$.\smallbreak
  If $n=1$, the result is trivial: $B(x,y)=axy$ for some $a\in \F$ is clearly symmetric.\smallbreak
	Fix $n\in\N$ and assume that every non-zero symmetric bilinear form on a dimension $n$ vector space over $\F$ is diagonalizable. Let $\dim V=n+1$. By the discussion below, $\exists \vx\in V$ such that $B(\vx,\vx)\neq 0$. Consider the linear map
\[\rT:V\to\F:\vv\mapsto B(\vx,\vv)\]
Clearly $\Rank\rT=1\implies \dim\cN(\rT)=n$. Moreover, $B$ is symmetric when restricted to $\cN(\rT)$; by the induction hypothesis there exists a basis $\beta$ of $\cN(\rT)$ such that $[\at{B}{\cN(\rT)}]_\beta$ is diagonal. But then $B$ is diagonal with respect to the basis $\beta\cup\{\vx\}$.\qedhere 
\end{enumerate}
\end{proof}



\begin{aside}{}{}
{\bfseries Aside: Characteristic two fields}\quad This means $1+1=0$ in $\F$; this holds, for instance, in the field $\Z_2=\{0,1\}$ of remainders modulo 2. We now see the importance of $\operatorname{char}\F\neq 2$ to the above result.
\begin{itemize}
	\item The proof uses the existence of $\vx\in V$ such that $B(\vx,\vx)\neq 0$. If $B$ is non-zero, $\exists\vu,\vv$ such that $B(\vu,\vv)\neq 0$. If both $B(\vu,\vu)=0=B(\vv,\vv)$, then $\vx=\vu+\vv$ does the job whenever $\operatorname{char}\F\neq 2$:
	\[B(\vx,\vx)=B(\vu,\vv)+B(\vv,\vu)=2B(\vu,\vv)\neq 0\]

	\item To see that the requirement isn't idle, consider $B(\vx,\vy)=\vx^T\begin{smatrix}0&1\\1&0\end{smatrix}\vy$ on the finite vector space $\Z_2^2=\{\stwovec 00,\stwovec 10,\stwovec 01,\stwovec 11\}$ over $\Z_2$. Every element of this space satisfies $B(\vx,\vx)=0$! Perhaps surprisingly, the matrix of $B$ is identical with respect to \emph{any} basis of $\Z_2^2$, whence $B$ is symmetric but non-diagonalizable.
\end{itemize}

% : choose any  and if both $B(\vu,\vu)=0=B(\vv,\vv)$, then $\vx=\vu+\vv$ satisfies
% \[B(\vx,\vx)=2B(\vu,\vv)\neq 0\]
% This is where the characteristic $\neq 2$ condition comes in!
% Why the problem with characteristic 2? Consider $B=\vv^T\begin{smatrix}0&1\\1&0\end{smatrix}\vw$. If we apply our algorithm:
% \begin{align*}
% &E_{12}^{(1)}\begin{pmatrix}
% 	0&1\\1&0
% \end{pmatrix}E_{21}^{(1)}=\begin{pmatrix}
% 	2&1\\1&0
% \end{pmatrix} \tag{add the second row to the first, columns similarly}\\
% &E_{21}^{(-1/2)}E_{12}^{(1)}\begin{pmatrix}
% 	0&1\\1&0
% \end{pmatrix}E_{21}^{(1)}E_{12}^{(-1/2)}=\begin{pmatrix}
% 	2&0\\0&-\frac 12
% \end{pmatrix} \tag{subtract $\frac 12$ row 1 from row 2, etc.}
% \end{align*}
% Since $Q=E_{21}^{(1)}E_{12}^{(-1/2)}= \begin{smatrix}
% 	1&0\\1&1
% \end{smatrix}\begin{smatrix}
% 	1&-\frac 12\\0&1
% \end{smatrix}=\begin{smatrix}
% 	1&-\frac 12\\1&\frac 12
% \end{smatrix}$
% we see that $B$ is diagonal with respect to the basis $\gamma=\left\{\stwovec 11,\stwovec{-1/2}{1/2}\right\}$. If the field has characteristic 2, however, then $\frac 12$ is meaningless! Perhaps surprisingly, the matrix of $B$ is the same with respect to \emph{any} basis of $\Z_2^2$: $B$ is therefore non-diagonalizable!
\end{aside}

\goodbreak


In Example \ref{ex:diagbilinear} notice how the three diagonal matrix representations have something in common: one each of the diagonal entries are positive and negative. This is a general phenomenon:

\begin{thm}{Sylvester's Law of Inertia}{}
Suppose $B$ is a symmetric bilinear form on a real vector space $V$ with diagonal matrix representation $\diag(\lambda_1,\ldots,\lambda_n)$. Then the number of entries $\lambda_j$ which are positive/negative/zero is independent of the diagonal representation.
\end{thm}

\begin{defn}{}{}
The \emph{signature} of a symmetric bilinear form $B$ is the triple $(n_+,n_-,n_0)$ representing how many positive, negative and zero terms are in any diagonal representation. Sylvester's Law says that the signature is an \emph{invariant} of a symmetric bilinear form.
\end{defn}

Positive-definiteness says that a real inner product on an $n$-dimensional space has signature $(n,0,0)$. Practitioners of relativity often work in \emph{Minkowski spacetime}: $\R^4$ equipped with a signature $(1,3,0)$ bilinear form, typically
\[B(\vx,\vy)=\vx^T\begin{pmatrix}
c^2&0&0&0\\
0&-1&0&0\\
0&0&-1&0\\
0&0&0&-1
\end{pmatrix}\vy =c^2x_1y_1-x_2y_2-x_3y_3-x_4y_4\]
where $c$ is the speed of light. Vectors are \emph{time-, space-,} or \emph{light-like} depending on whether $B(\vx,\vx)$ is positive, negative or zero. For instance $\vx=3c^{-1}\ve_1+2\ve_2+2\ve_3+\ve_4$ is light-like.

\begin{proof}
For simplicity, let $V=\R^n$ and write $B(\vx,\vy)=\vx^TA\vy$ where $A$ is symmetric.
\begin{enumerate}
  \item Define $\Rank B:=\Rank A$ and observe this is independent of basis (exercises).
  \item Let $\beta,\gamma$ be diagonalizing bases, ordered according to whether $B$ is \textcolor{red}{positive}, \textcolor{blue}{negative} or \textcolor{Green}{zero}.
  \begin{gather*}
  \beta=\{\textcolor{red}{\vv_1},\ldots,\textcolor{red}{\vv_p}, \textcolor{blue}{\vv_{p+1}},\ldots,\textcolor{blue}{\vv_r}, \textcolor{Green}{\vv_{r+1}},\ldots,\textcolor{Green}{\vv_n}  \}\\
  \gamma=\{\textcolor{red}{\vw_1},\ldots,\textcolor{red}{\vw_q}, \textcolor{blue}{\vw_{q+1}},\ldots,\textcolor{blue}{\vw_r}, \textcolor{Green}{\vw_{r+1}},\ldots,\textcolor{Green}{\vw_n}  \}
  \end{gather*}
  Here $r=\Rank B$ in accordance with part 1: our goal is to prove that $p=q$.
  \item Assume $p<q$, define the matrix $C$ and check what follows:
  \[C=\begin{smatrix}
    \textcolor{red}{\vv_1^T}A\\\vdots\\\textcolor{red}{\vv_p^T}A\\\textcolor{blue}{\vw_{q+1}^T}A\\\vdots\\\textcolor{blue}{\vw_{r}^T}A
    \end{smatrix}\in M_{(r-q+p)\times n}(\R)\]
    \begin{enumerate}
      \item $\Rank C\le r-q+p<r\implies \Null C> n-r$, thus
      \[\exists\vx\in\R^n\text{ such that }C\vx=\V0 \text{ and }\vx\not\in\Span\{\textcolor{Green}{\vv_{r+1}},\ldots,\textcolor{Green}{\vv_n}\}\]
      \item The first $p$ entries of $C\vx=\V0$ mean that $\vx\in\Span\{\textcolor{blue}{\vv_{p+1}},\ldots,\textcolor{blue}{\vv_r}, \textcolor{Green}{\vv_{r+1}},\ldots,\textcolor{Green}{\vv_n}\}$ and so $B(\vx,\vx)<0$. Note how we use part (a) to get a strict inequality here.
      \item Now write $\vx$ with respect to $\gamma$: this time we see that $\vx\in\Span\{\textcolor{red}{\vw_1},\ldots,\textcolor{red}{\vw_q}, \textcolor{Green}{\vw_{r+1}},\ldots,\textcolor{Green}{\vw_n}\}$, whence $B(\vx,\vx)\ge 0$: a contradiction.\qedhere
 		\end{enumerate}
\end{enumerate}
\end{proof}

\goodbreak

\boldsubsubsection{Quadratic Forms \& Diagonalizing Conics}


\begin{defn}{}{}
To every symmetric bilinear form $B:V\times V\to\F$ is associated a \emph{quadratic form}
\[K:V\to\F:\vx\mapsto B(\vx,\vx)\]
A function $K:V\to\F$ is termed a quadratic form when such a symmetric bilinear form exists.
\end{defn}


\begin{examples}{}{}
\exstart If $B$ is a real inner product, then $K(\vv)=\ip{\vv,\vv}=\Nm{\vv}^2$ is the square of the norm.
\begin{enumerate}\setcounter{enumi}{1}  
  \item Let $\dim V=n$ and $A$ be the matrix of $B$ with respect to a basis $\beta$. By the symmetry of $A$,
  \[K(\vx)=\vx^TA\vx=\sum_{i,j=1}^nx_iA_{ij}x_j =\sum_{1\le i\le j\le n}\tilde a_{ij}x_ix_j
  \quad\text{where}\ \tilde a_{ij}=\begin{cases}
  A_{ij}&\text{if $i=j$}\\
  2A_{ij}&\text{if }i\neq j
  \end{cases}\]
  E.g., $K(\vx)=3x_1^2+4x_2^2-2x_1x_2$ corresponds to the bilinear form $B(\vx,\vy)=\vx^T\begin{smatrix}
  3&-1\\-1&4
  \end{smatrix}\vy$
%   where $\tilde a_{ij}=\begin{cases}
%   A_{ij}&\text{if $i=j$}\\
%   2A_{ij}&\text{if }i\neq j
%   \end{cases}$
\end{enumerate}
\end{examples}


As a fun application, we consider the diagonalization of conics in $\R^2$. The general non-zero conic has equation
\[ax^2+2bxy+cy^2+dx+ey+f=0\]
where the first three terms comprise a quadratic form
\[K(\vx)=ax^2+2bxy+cy^2\leftrightsquigarrow B(\vv,\vw)=\vv^T\begin{pmatrix}
a&b\\b&c
\end{pmatrix}\vw\]
If $\{\vv_1,\vv_2\}$ is a diagonalizing basis, then there exist scalars $\lambda_1,\lambda_2$ for which
\[K(t_1\vv_1+t_2\vv_2)=\lambda_1t_1^2+\lambda_2t_2^2\]
whence the general conic becomes
\[\lambda_1t_1^2+\lambda_2t_2^2+\mu_1t_1+\mu_2t_2=\eta,\quad \lambda_1,\lambda_2,\mu_1,\mu_2,\eta\in\R\]
If $\lambda_i\neq 0$, we may complete the square via the linear transformation $s_j=t_j+\frac{\mu_j}{2\lambda_j}$. The canonical forms are then recovered:
\begin{itemize}
	\item[]\normalfont\emph{Parabola}: $\lambda_1$ or $\lambda_2=0$ (but not both).
	\item[]\normalfont\emph{Ellipse}: $\lambda_1s_1^2+\lambda_2s_2^2=k\neq 0$ where $\lambda_1,\lambda_2,k$ have the same sign.
	\item[]\normalfont\emph{Hyperbola}: $\lambda_1s_1^2+\lambda_2s_2^2=k\neq 0$ where $\lambda_1,\lambda_2$ have opposite signs.
\end{itemize}
Since $B$ is symmetric, we may take $\{\vv_1,\vv_2\}$ to be an \emph{orthonormal} basis of $\R^2$, whence any conic may be put in canonical form by applying only a rotation/reflection and translation (completing the square).\smallbreak
Alternatively, we could diagonalize $K$ using our earlier algorithm; this additionally permits shear transforms. By Sylvester's Law, the diagonal entries will have the same number of $(+,-,0)$ terms regardless of the method, so the canonical form will be unchanged.


\begin{examples}{}{}
\exstart We describe and plot the conic with equation $7x^2+24xy=144$.
\begin{enumerate}\setcounter{enumi}{1} 
\begin{minipage}[t]{0.6\linewidth}\vspace{-10pt}
  \item[] The matrix of the associated bilinear form is $\begin{smatrix}
7&12\\12&0
  \end{smatrix}$ which has orthonormal eigenbasis
	\[\beta=\{\vv_1,\vv_2\}=\left\{\frac 15\twovec 43,\frac 15\twovec{-3}4\right\}\]
	with eigenvalues $(\lambda_1,\lambda_2)=(16,-9)$. In the rotated basis, this is the canonical hyperbola
	\[16t_1^2-9t_2^2=144\iff \frac{t_1^2}{3^2}-\frac{t_2^2}{4^2}=1\]
	which is easily plotted. In case this is too fast, use the change of co-ordinate matrix to compute directly:
\end{minipage}\begin{minipage}[t]{0.4\linewidth}\vspace{-5pt}
\flushright\includegraphics[scale=0.9]{quad-hyperbola}
\end{minipage}\smallbreak
\vspace{-10pt}
	\[Q_\beta^\epsilon =\frac 15\begin{pmatrix}
	4&-3\\3&4
	\end{pmatrix}\implies \twovec{t_1}{t_2}=[\vx]_\beta=Q_\epsilon^\beta\twovec xy=
	\frac 15\twovec{4x+3y}{-3x+4y}\]
	which quickly recovers the original conic by substitution.


\begin{minipage}[t]{0.64\linewidth}\vspace{0pt}
\item The conic defined by $K(\vx)=x^2+12xy+3y^2=-33$ defines a hyperbola in accordance with Example \ref{ex:diagbilinear}. With respect to the basis $\gamma=\{\stwovec 1{-2},\stwovec 01\}$, we see that
\begin{gather*}
K(\vx)=-11x^2+3(2x+y)^2=-11t_1^2+3t_2^2=-33\\
\iff \frac{t_1^2}{3}-\frac{t_2^2}{11}=1
\end{gather*}
Even though $\gamma$ is non-orthogonal, we can still plot the conic! If we instead used the orthonormal basis $\eta$, we'd obtain
\[K(\vx)=(\sqrt{37}+2)s_1^2-(\sqrt{37}-2)s_2^2=-33\]
however the calculation to find $\eta$ is time-consuming and the expressions for $s_1,s_2$ are extremely ugly.
\end{minipage}\begin{minipage}[t]{0.36\linewidth}\vspace{0pt}
\flushright\includegraphics[scale=0.97]{quad-hyperbola2}
\end{minipage}
\end{enumerate}
\end{examples}

A similar approach can be applied to higher degree quadratic equations/manifolds: e.g.\ ellipsoids, paraboloids and hyperboloids in $\R^3$.

\vfil
 
\begin{exercises}
\exstart Prove that the sum of any two bilinear forms is bilinear, and that any scalar multiple of a bilinear form is bilinear: thus the set of bilinear forms on $V$ is a vector space.\vspace{-5pt}
  
\begin{enumerate}\setcounter{enumi}{1}
  \item[](\emph{You can't use matrices here, since $V$ could be infinite-dimensional!}) 
  \item Compute the matrix of the bilinear form
  \[B(\vx,\vy)=x_1y_1-2x_1y_2+x_2y_1-x_3y_3\]
  on $\R^3$ with respect to the basis $\beta=\left\{\sthreevec 101,\sthreevec 10{-1},\sthreevec 010\right\}$.
  
  \goodbreak
  
  \item Check that the function $B(f,g)=f'(0)g''(0)$ is a bilinear form on the vector space of twice-differentiable functions. Find the matrix of $B$ with respect to $\beta=\{\cos t,\sin t,\cos 2t,\sin 2t\}$ when restricted to the subspace $\Span\beta$.
  
  \item For each matrix $A$, find a diagonal matrix $D$ and an invertible matrix $Q$ such that $Q^TAQ=D$.
  \begin{enumerate}
    \item $\begin{pmatrix}
    1&3\\3&2
    \end{pmatrix}$\qquad
    (b)\ \ $\begin{pmatrix}
    3&1&2\\1&4&0\\2&0&-1
    \end{pmatrix}$
  \end{enumerate}
  
  \item If $K$ is a quadratic form and $K(\vx)=2$, what is the value of $K(3\vx)$?
  
  \item If $\F$ does not have characteristic 2, and $K(\vx)=B(\vx,\vx)$ is a quadratic form, prove that we can recover the bilinear form $B$ via
		\[B(\vx,\vy)=\frac 12\left(K(\vx+\vy)-K(\vx)-K(\vy)\right)\]
  
  \item If $B(\vx,\vy)=\vx^T\begin{smatrix}0&1\\1&0\end{smatrix}\vy$ is a bilinear form on $\F^2$, compute the quadratic form $K(\vx)$.
  
  \item Suppose $B$ is a symmetric bilinear form on a real finite-dimensional space. With reference to the proof of Sylvester's Law, explain why $\Rank B$ is independent of the choice of diagonalizing basis.
  
  \item If $\operatorname{char}\F\neq 2$, apply the diagonalizing algorithm to the symmetric bilinear form $B=\vx^T\begin{smatrix}0&1\\1&0\end{smatrix}\vy$ on $\F^2$. What goes wrong if $\operatorname{char}\F=2$? 
  
  \item Describe and plot the following conics:
  \begin{enumerate}
    \item $x^2+y^2+xy=6$
    
    \item $35x^2+120xy=4x+3y$
  \end{enumerate}
  
  \item Suppose that a non-empty, non-degenerate\footnotemark conic $C$ in $\R^2$ has the form $ax^2+2bxy+cy^2+dx+ey+f=0$, where at least one of $a,b,c\neq 0$, and define $\Delta=b^2-ac$. Prove that:
  \begin{itemize}
    \item $C$ is a parabola if and only if $\Delta=0$;
    \item $C$ is an ellipse if and only if $\Delta<0$;
    \item $C$ is a hyperbola if and only if $\Delta>0$.
  \end{itemize}
  (\emph{Hint: $\lambda_1,\lambda_2$ are the eigenvalues of a symmetric matrix, so\ldots})
  
  %\item Consider the quadratic form $K(x,y,z)=x^2+y^2-z^2$ which describes a cone in three dimensions. Let $ax+by+cz=1$ be the equation of a plane in $\R^3$ which does not pass through the origin. Prove that 
\end{enumerate}
\end{exercises}

\footnotetext{The conic contains at least two points and cannot be factorized as a product of two straight lines: for example, the following are disallowed;
  \begin{itemize}
    \item $x^2+y^2+1=0$ is empty (unless one allows conics over $\C\ldots$)
    \item $x^2+y^2=0$ contains only one point;
    \item $x^2-xy-x+y=(x-1)(x-y)=0$ is the product of two lines.
  \end{itemize}}
\clearpage

\iffalse

\setcounter{subsection}{10}
\subsection{The Geometry of Orthogonal Operators}

The purpose of this section is to classify the orthogonal operators $T\in\cL(V)$ on any finite-dimensional \emph{real} inner product space $V$. This is very easy to do when $\dim V=1$ or 2.

\begin{description}
	\item[\normalfont\emph{Dimension 1}:] Clearly $T(\vx)=a\vx$ for some $a\in\R$. $T$ is clearly self-adjoint and therefore orthogonal if and only if $a^2=1$, yielding $T=\pm I$.
	\item[\normalfont\emph{Dimension 2}:] Any choice of orthonormal basis $\beta$ generates an isomorphism of inner product spaces $\phi:V\to\R^2$ and the problem becomes that of identifying every orthogonal matrix in $M_2(\F)$. We already did this in Exercise \ref{ex:rotrefmatrix}: all orthogonal $2\times 2$ real matrices have the form
\[A_\theta=\begin{pmatrix}
	\cos\theta&-\sin\theta\\
	\sin\theta&\cos\theta
  \end{pmatrix},\qquad B_\theta=\begin{pmatrix}
	\cos\theta&\sin\theta\\
	\sin\theta&-\cos\theta
	\end{pmatrix}\]
for some $\theta\in\R$. Moreover, their effects are as follows:
	\begin{itemize}
	  \item $L_{A_\theta}\in\cL(\R^2)$ rotates vectors counter-clockwise about the origin by $\theta$ radians;
	  \item $L_{B_\theta}\in\cL(\R^2)$ reflects vectors across the line making angle $\frac 12\theta$ with the positive $x$-axis.
	\end{itemize}
\end{description}
	
These maps provides the motivation for a general definition of rotation and reflection.
	
\begin{defn}{}{}
Suppose $T\in\cL(V)$ where $V$ is a finite-dimensional real inner product space. We say that $T$ is a:
\begin{description}
  \item[\normalfont\emph{Rotation}] if either $T$ is the identity map, or if there exists a two-dimensional subspace $W\le V$ with orthonormal basis $\beta=\{\vx,\vy\}$ and some $\theta\in(0,2\pi)$ such that\footnote{We could include the identity here with $\theta=0$, but then \emph{any} 2-dimensional subspace $W$ is suitable and the axis of rotation is not well-defined.}
  \[T(\vx)=\cos\theta\,\vx+\sin\theta\,\vy,\qquad T(\vy)=-\sin\theta\,\vx+\cos\theta\,\vy,\qquad T(\vw^\perp)=\vw^\perp,\]
  for all $\vw^\perp\in W^\perp$, termed the \emph{axis of rotation.} Otherwise said, $[\at T{W}]_\beta=A_\theta$ and $\at T{W^\perp}=I$. 
  
  \item[\normalfont\emph{Reflection}] if there exists a one-dimensional subspace $W\le V$ such that $\forall \vw\in W,\vw^\perp\in W^\perp$,
  \[T(\vw+\vw^\perp)=-\vw+\vw^\perp\]
  Otherwise said, $\at T{W}=-I$ and $\at T{W^\perp}=I$. We could call this \emph{reflection across $W^\perp$.}
\end{description}
\end{defn}

\begin{example}{}{}
We compute the matrix of the reflection across $W^\perp=\Span\{\vi+2\vj\}\le\R^2$.\\[5pt]
First observe that $W=\Span\{2\vi-\vj\}$, whence $\beta=\{\stwovec 2{-1},\stwovec 12\}$ is an orthogonal basis (it doesn't need to be orthonormal!). The change of co-ordinate matrix from $\beta$ to the standard basis $\epsilon=\{\vi,\vj\}$ is $Q=Q_\beta^\epsilon=\begin{smatrix}
2&1\\-1&2
\end{smatrix}$, whence the required matrix is
\[A=Q_\beta^\epsilon\begin{pmatrix}
-1&0\\0&1
\end{pmatrix}Q_\epsilon^\beta =\frac 15\begin{pmatrix}
-3&4\\4&3
\end{pmatrix}\]
which is easily checked to satisfy $A\stwovec 2{-1}=-\stwovec 2{-1}$ and $A\stwovec 12=\stwovec 12$.
\end{example}

It should be clear that every rotation and reflection is an orthogonal operator and that they may be easily distinguished by determinant:
\[\det(\text{rotation})=1,\qquad \det(\text{reflection})=-1\]
It is also easy to see that $L_{A_\theta}\in\cL(\R^2)$ is a rotation as per the definition: simply take $\beta=\{\vi,\vj\}$ and $W=\R^2$! We leave the verification that $L_{B_\theta}\in\cL(\R^2)$ is a reflection as an exercise. The case for $\dim V=1$ is even easier. To summarize:

\begin{lemm}{}{22orth}
Suppose $V$ is a real inner product space with $\dim V=1$ or $2$. Then every orthogonal operator $T\in\cL(V)$ is either a rotation or a reflection.
\end{lemm}


Our goal is to describe any orthogonal operator $T$ in terms of rotations and reflections. The upshot is that there exists an orthonormal basis $\beta$ with respect to which the matrix of $T$ has block form
\[[T]_\beta=\diag(R_1,R_2,\ldots,R_k)\]
and each $R_j$ is either a $2\times 2$ rotation matrix, or is $\pm 1$.

\begin{example}{}{}
The matrix $A=\begin{smatrix}
\frac 35&-\frac 45&0&0\\
\frac 45&\frac 35&0&0\\
0&0&1&0\\
0&0&0&-1
\end{smatrix}$ describes an orthogonal operator $L_A\in\cL(\R^4)$. We have three commuting orthogonal operations $L_A=T_1T_2T_3$:
\begin{itemize}
  \item $T_1$ rotates by $\theta=\tan^{-1}\frac 43$ about the axis $W_1^\perp=\Span\{\ve_3,\ve_4\}$;
  \item $T_2$ is the identity rotation with `axis' $W_2^\perp=\Span\{\ve_1,\ve_2,\ve_4\}$;
  \item $T_3$ is reflection across $W_3^\perp=\Span\{\ve_1,\ve_2,\ve_3\}$
\end{itemize}
\end{example}

To see that this always happens, we first need a little theory regarding a real $n\times n$ matrix $A$. Our primary strategy relies on eigendecompositions; unfortunately $A$ need not have any \emph{real} eigenvectors. As a linear operator on $\C^n$ however, it must have at least one eigenpair $(\lambda,\vv)$: consider the real and imaginary parts: $\vv=\vx+i\vy$ and $\lambda=p+iq$ and observe that
\[A\vx+iA\vy=A(\vx+i\vy)=(p+iq)(\vx+i\vy)=(p\vx-q\vy)+i(p\vy+q\vx) \tag{$\ast$}\]
Since $A$ is real, we can equate real and imaginary parts to see that $\Span\{\vx,\vy\}$ \emph{is an $L_A$-invariant subspace of $\R^n$.} This is the generic situation, since, upon choosing a basis, any linear operator on a finite-dimensional real vector space is multiplication by some real matrix $A$. We therefore conclude:

\begin{lemm}{}{}
If $T\in\cL(V)$ is a linear operator on a finite-dimensional real vector space, then there exists a $T$-invariant subspace $W\le V$ such that $\dim W=1$ or 2.
\end{lemm}

The reason our construction might product a 1-dimensional space? The real and imaginary parts $\vx,\vy$ might be parallel! In this case, of course, $W$ is an eigenspace of $T$.\\

Of particular interest to our current goal is that the real and imaginary parts of ($\ast$) look very like a \emph{rotation.} Indeed if we specialize the Lemma to an orthogonal operator $T$, it is an easy exercise to check that $W^\perp$ is also $T$-invariant. We conclude that $\at T{W^\perp}$ is an orthogonal operator on a lower-dimensional vector space. We may now repeat the process until we have exhausted $V$. This leads to the following:

\begin{thm}{}{}
Let $T\in\cL(V)$ be an orthogonal operator on a finite-dimensional real inner product space. Then:
\begin{enumerate}
  \item There exists an orthogonal decomposition
	\[V= W_1\oplus\cdots \oplus W_m\]
	where each $W_k$ is a 1 or 2-dimensional $T$-invariant subspace.\footnote{It should be clear that $\frac n2\le m\le n$.}
	\item For each $j$, define $T_j:=\at T{W_j}+\at I{W_j^\perp}$. Then $T_j$ is a rotation or a reflection and $T=T_1\cdots T_m$.
	\item $\det T=(-1)^k$ where exactly $k$ of the $T_j$ are reflections.
	\item The subspaces may be chosen so that there is at most one reflection $T_j$, and where $\dim W_j=1$.
\end{enumerate}
\end{thm}

\begin{proof}
\begin{enumerate}
  \item We've already done this, though it could be made more formal by induction.
  \item Since $\at T{W_j}$ is orthogonal on a 1 or 2-dimensional space, Lemma \ref{lemm:22orth} says it is either a rotation or a reflection. Extending to $T_j$ by forcing the identity on $W_j^\perp$ produces a rotation or reflection on $\R^n$. Moreover, since the $W_j$ are orthogonal subspaces, we see that
  \[\vv_j\in W_j\implies T_k(\vv_j)=\begin{cases}
  T(\vv_j)&\text{if $j=k$}\\
  \vv_j&\text{if $j\neq k$}
  \end{cases}\]
  By applying to a basis where each $\vv_j$ lies in some $W_j$, the relationship $T=T_1\cdots T_m$ is now clear.
  \item This is immediate since every rotation has determinant 1 and every reflection determinant $-1$.
  \item Suppose that $T_j$ is a reflection and $\dim W_j=2$. Then there exists a one-dimensional subspace $\hat W_j$ such that
  \[\at{{T_j}}{{\hat W_j}}=-I,\qquad \at{{T_j}}{{\hat W_j^\perp}}=I\]
  By replacing $W_j$ by $\hat W_j\oplus \hat W_j^\perp$, we see that every reflection may be assumed to be one-dimensional.\\
  Now suppose $W_j,W_k$ are 1-dimensional subspaces corresponding to reflections $T_j,T_k$. Then
  \[(T_j+T_k)(\vv)=\begin{cases}
  -\vv&\text{if $\vv\in W_j\oplus W_k$}\\
  \vv&\text{if $\vv\in (W_j\oplus W_k)^\perp$}
  \end{cases}\]
  Otherwise said, $T_j+T_k$ is the \emph{rotation} with $W=W_j\oplus W_k$ and $\theta=\pi$. We can therefore replace an even number of reflections with rotations.\qedhere
\end{enumerate}
\end{proof}



% 
% \begin{lemm}{}{orthrotref}
% Every rotation is an orthogonal operator with determinant 1.\\
% Every reflection is an orthogonal operator with determinant $-1$.
% \end{lemm}\pagebreak[3]
% 
% \begin{proof}
% Given a rotation $T\in\cL(V)$ with axis $W^\perp$, extend the basis $\beta=\{\vx,\vy\}$ of $W$ to an orthonormal basis $\gamma$ of $V$. With respect to this basis, the matrix of $T$ has block form
% \[[T]_\gamma=\left(\begin{array}{c|c}
% A_\theta&0\\\hline
% 0&I
% \end{array}\right) \implies [T^*]_\gamma=\left(\begin{array}{c|c}
% A_\theta^T&0\\\hline
% 0&I
% \end{array}\right)
% =\left(\begin{array}{c|c}
% A_\theta^{-1}&0\\\hline
% 0&I
% \end{array}\right) =[T]_\gamma^{-1}\]
% since $A_\theta$ is orthogonal. We conclude that $T$ is invertible with inverse $T^{-1}=T^*$: that is, $T$ is an orthogonal operator. Moreover its determinant is clearly $\det A_\theta=1$.\\[5pt]
% We leave the reflection argument as an exercise.
% \end{proof}



\begin{example}{}{orthorotatenasty}
Consider the matrix $A=\frac 1{25}\begin{smatrix}
  15&-12&-16\\
  20&9&12\\
  0&-20&15
  \end{smatrix}$. We show that $L_A\in\cL(\R^3)$ is a rotation and find its axis and angle of rotation.\\[5pt]
  It is easy (with a calculator!) to check that $A^TA=I$ so that $A$ is orthogonal, and moreover that $\det A=1$.
  %\[\det A=\frac 1{25^3}\left(15^2\cdot 9+16\cdot 20^2+2\cdot 12\cdot 15\cdot 20\right)=1\]
  According to the Theorem, there are no reflections in $A$ and so there must exist a decomposition $\R^3=W\oplus W^\perp$ where $\dim W=2$ and such that $W^\perp$ is an eigenspace with eigenvalue 1. We can hunt for the eigenvalues via the usual approach, or we can do it by brute force: $\vw^\perp=\sthreevec abc$ has eigenvalue 1 if and only if
   \[\begin{cases}
   15a-12b-16c=25a\\
   20a+9b+12c=25b\\
   -20b+15c=25c
   \end{cases} \iff c=-2b, \ a=2b \iff \vw^\perp\propto\threevec 21{-2}\]
   The axis of rotation is therefore $W^\perp=\Span\left\{\sthreevec 21{-2}\right\}$. To find the angle, we simply need to choose anything orthogonal to this (and therefore in $W$), and see how $A$ changes the angle. Thus let $\vw=\sthreevec 101$ and compute
   \[A\vw=\frac 1{25}\threevec{-1}{32}{15}\implies \cos\theta=\frac{\sthreevec 101\cdot\sthreevec{-1}{32}{15}}{ \sqrt{1+1}\sqrt{1+32^2+15^2}}
   =\frac{14}{50}=\frac 7{25} \implies \theta\approx \ang{74}\]
   In fact $(\vw,A\vw,\vw^\perp)$ is negatively oriented if we `look down' on $W=\Span\{\vw,A\vw\}$ from above $\vw^\perp$, the rotation would be through  \ang{74} \emph{clockwise.}
\end{example}

Because of the theorem, some authors prefer to simply define a rotation to be any orthogonal operator with determinant 1 and a reflection to be any with determinant $-1$. The multiplicative property of determinant then shows that the composition of any two rotations is a rotation,\footnote{In particular, this shows that the set of `rotations' forms a normal subgroup $\rSO(V)$ of the group of orthogonal operators on $V$, known as the \emph{special orthogonal group.}} of two reflections is a rotation, etc. The problem with this definition is that one loses the notion of an axis of rotation, since a general determinant 1 orthogonal operator will typically require multiple rotations around several axes.

\begin{example}{}{}
For given $\theta,\phi$, the matrix
\[A=\begin{smatrix}
  \cos\theta&-\sin\theta&0&0\\
  \sin\theta&\cos\theta&0&0\\
  0&0&\cos\phi&-\sin\phi\\
  0&0&\sin\phi&\cos\phi
\end{smatrix}\]
performs two rotations. It has no `axis' of rotation, since it has no eigenspace with eigenvalue 1: indeed it has no \emph{real} eigenvalues at all. Regardless, it is orthogonal and has determinant 1, so many authors consider this to be a rotation.
\end{example}

\begin{exercises}
\exstart Find the matrices with respect to the standard basis of the reflections in $W^\perp$:

\begin{enumerate}
	\item[]\begin{enumerate}
	  \item $W=\Span\{-2\vi+5\vj\}$
	  \item $W=\Span\{\vi+2\vj+2\vk\}$\quad  \emph{Hint: consider the basis $\beta=\left\{\frac 13\sthreevec 122,\frac 13\sthreevec 2{-2}1,\frac 13\sthreevec 21{-2}\right\}$}
	\end{enumerate}
	
  \item\begin{enumerate}
  	\item Given $\phi$, let $\beta=\{\vx,\vy\}=\left\{\stwovec{\cos\phi}{\sin\phi},\stwovec{-\sin\phi}{\cos\phi}\right\}$ and $\gamma=\{\vx,-\vy\}$. Clearly every orthonormal basis of $\R^2$ has the form $\beta$ or $\gamma$. Compute the matrices
  	\[[L_{A_\theta}]_\beta\quad\text{and}\quad[L_{A_\theta}]_\gamma\quad\text{where}\quad A_\theta=\begin{pmatrix}
		\cos\theta&-\sin\theta\\
		\sin\theta&\cos\theta
  	\end{pmatrix}\]
  	What do you observe?
  	\item Now consider the reflection matrix $B_\theta=\begin{smatrix}
		\cos\theta&\sin\theta\\
		\sin\theta&-\cos\theta
  	\end{smatrix}$. Show that there is a privileged orthonormal basis (indeed four up to scaling and reordering) $\beta$ such that $[L_{B_\theta}]_\beta=\begin{smatrix}
		1&0\\
		0&-1
  	\end{smatrix}$.
 	\end{enumerate}
 	
 	\item Compute $A_\phi B_\theta A_\phi^{-1}$. Across what line in $\R^2$ does this matrix reflect? Can you do this \emph{without} multiplying out the matrices?
  
  \item Prove explicitly:
  \begin{enumerate}
    \item If $T$ is a rotation then it is orthogonal and $\det T=1$;
    \item If $T$ is a reflection then it is orthogonal and $\det T=-1$.
  \end{enumerate}
  
  \item If $W$ is a $T$-invariant subspace where $T$ is an orthogonal operator on a finite-dimensional real inner product space $V$, prove that $W^\perp$ is also $T$-invariant.

	\item Find the result when the vector $\vx=\sthreevec 342$ is rotated \ang{90} counter-clockwise around the $z$-axis (when viewed from above) and then reflected in the $xz$-plane. What is the effect of this transformation on a general vector?
  
  \item For any non-zero $\theta,\phi\in(-\pi,\pi)$, define the matrices
  \[A=\begin{pmatrix}
  \cos\theta&-\sin\theta&0\\
  \sin\theta&\cos\theta&0\\
  0&0&1
  \end{pmatrix}\qquad B=\begin{pmatrix}
  1&0&0\\
  0&\cos\phi&-\sin\phi\\
  0&\sin\phi&\cos\phi
  \end{pmatrix}\]
  Prove that $L_{AB}$ is a rotation of $\R^3$ and find its axis of rotation.

 	
%  	 since it has eigenspaces
% \[W=E_{-1}=\Span\left\{\twovec{-\sin\frac\theta 2}{\cos\frac\theta 2}\right\},\qquad W^\perp=E_1=\Span\left\{\twovec{\cos\frac\theta 2}{\sin\frac\theta 2}\right\}\]
  
%  
%   
%   \item Recall (Exercise \ref{ex:unitarymatrix}) that a unitary matrix $A\in M_n(\C)$ has an orthonormal eigenbasis $\gamma=\{\vect{\vv}{n}\}$ and eigenvalues $e^{i\theta_1},\ldots,e^{i\theta_n}$.
%   \begin{enumerate}
%     \item Suppose $A$ is real (it is an \emph{orthogonal} matrix). Prove that each $\cl{\vv_j}$ is also an eigenvector with eigenvalue $e^{-i\theta_j}$.
%     \item If $e^{i\theta_j}$ is not real, show that $\dim\Span\{\vv_j,\cl{\vv_j}\}=2$. Moreover, find a 2-dimensional subspace $W$ of $\R^n$ so that $L_A$ restricted to $W$ is a rotation by angle $\theta_j$.
%     \item If $e^{i\theta_j}$ is real, show that there exists a 1 or 2-dimensional subspace $W$ of $\R^n$ such that $L_A$ restricted to $W$ is $\pm I$.
% 	\end{enumerate}
	
	\item Let $V$ be a real inner product space of dimension 2. Given any $\vx\neq \vy$ with $\Nm\vx=\Nm\vy=1$, show that there exists a unique rotation $T\in\cL(V)$ such that $T(\vx)=\vy$.
% 	  \item (If you've done group actions in group theory) Write $\rSO(n+1)$ for the special orthogonal group $\rSO(\R^{n+1})$. Part (a) shows that the action of the group $\rSO(n+1)$ is transitive on $S^n$. Given $\vx\in S^n$ consider the \emph{stabilizer} subgroup
% 	  \[\Stab(\vx)=\{U\in\rSO(n+1):U(\vx)=\vx\}\]
% 	  Prove that $\Stab(\vx)\cong\rSO(n)$ and moreover that there exists a bijection of the unit sphere with the set of left cosets
% 	  \[\mu:S^n\to\quotient{\rSO(n+1)}{\rSO(n)}=\{TU:T\in\rSO(n+1),U\in\Stab(\vx)\}\]
% 	  (\emph{This `transitive action modulo stabilizer' construction is known as a homogeneous space})
% 	\end{enumerate}
\end{enumerate}
\end{exercises}

\fi